\documentclass{report}

%%%%%%%%%%%%%%%%%%%%%%%%%%%%%%%%%
% PACKAGE IMPORTS
%%%%%%%%%%%%%%%%%%%%%%%%%%%%%%%%%


\usepackage[tmargin=2cm,rmargin=1in,lmargin=1in,margin=0.85in,bmargin=2cm,footskip=.2in]{geometry}
\usepackage{amsmath,amsfonts,amsthm,amssymb,mathtools}
\usepackage[varbb]{newpxmath}
\usepackage{xfrac}
\usepackage[makeroom]{cancel}
\usepackage{mathtools}
\usepackage{bookmark}
\usepackage{enumitem}
\usepackage{hyperref,theoremref}
\hypersetup{
	pdftitle={Assignment},
	colorlinks=true, linkcolor=doc!90,
	bookmarksnumbered=true,
	bookmarksopen=true
}
\usepackage[most,many,breakable]{tcolorbox}
\usepackage{xcolor}
\usepackage{varwidth}
\usepackage{varwidth}
\usepackage{etoolbox}
%\usepackage{authblk}
\usepackage{nameref}
\usepackage{multicol,array}
\usepackage{tikz-cd}
\usepackage[ruled,vlined,linesnumbered]{algorithm2e}
\usepackage{comment} % enables the use of multi-line comments (\ifx \fi) 
\usepackage{import}
\usepackage{xifthen}
\usepackage{pdfpages}
\usepackage{transparent}
\usepackage[utf8]{inputenc}
\usepackage{fontawesome}


\newcommand\mycommfont[1]{\footnotesize\ttfamily\textcolor{blue}{#1}}
\SetCommentSty{mycommfont}
\newcommand{\incfig}[1]{%
    \def\svgwidth{\columnwidth}
    \import{./figures/}{#1.pdf_tex}
}

\usepackage{tikzsymbols}
\renewcommand{\qedsymbol}{\faWhatsapp}


%\usepackage{import}
%\usepackage{xifthen}
%\usepackage{pdfpages}
%\usepackage{transparent}


%%%%%%%%%%%%%%%%%%%%%%%%%%%%%%
% SELF MADE COLORS
%%%%%%%%%%%%%%%%%%%%%%%%%%%%%%



\definecolor{myg}{RGB}{56, 140, 70}
\definecolor{myb}{RGB}{45, 111, 177}
\definecolor{myr}{RGB}{199, 68, 64}
\definecolor{mytheorembg}{HTML}{F2F2F9}
\definecolor{mytheoremfr}{HTML}{00007B}
\definecolor{mylenmabg}{HTML}{FFFAF8}
\definecolor{mylenmafr}{HTML}{983b0f}
\definecolor{mypropbg}{HTML}{f2fbfc}
\definecolor{mypropfr}{HTML}{191971}
\definecolor{myexamplebg}{HTML}{F2FBF8}
\definecolor{myexamplefr}{HTML}{88D6D1}
\definecolor{myexampleti}{HTML}{2A7F7F}
\definecolor{mydefinitbg}{HTML}{E5E5FF}
\definecolor{mydefinitfr}{HTML}{3F3FA3}
\definecolor{notesgreen}{RGB}{0,162,0}
\definecolor{myp}{RGB}{197, 92, 212}
\definecolor{mygr}{HTML}{2C3338}
\definecolor{myred}{RGB}{127,0,0}
\definecolor{myyellow}{RGB}{169,121,69}
\definecolor{myexercisebg}{HTML}{F2FBF8}
\definecolor{myexercisefg}{HTML}{88D6D1}


%%%%%%%%%%%%%%%%%%%%%%%%%%%%
% TCOLORBOX SETUPS
%%%%%%%%%%%%%%%%%%%%%%%%%%%%

\setlength{\parindent}{1cm}
%================================
% THEOREM BOX
%================================

\tcbuselibrary{theorems,skins,hooks}
\newtcbtheorem[number within=section]{Theorem}{Theorem}
{%
	enhanced,
	breakable,
	colback = mytheorembg,
	frame hidden,
	boxrule = 0sp,
	borderline west = {2pt}{0pt}{mytheoremfr},
	sharp corners,
	detach title,
	before upper = \tcbtitle\par\smallskip,
	coltitle = mytheoremfr,
	fonttitle = \bfseries\sffamily,
	description font = \mdseries,
	separator sign none,
	segmentation style={solid, mytheoremfr},
}
{th}

\tcbuselibrary{theorems,skins,hooks}
\newtcbtheorem[number within=chapter]{theorem}{Theorem}
{%
	enhanced,
	breakable,
	colback = mytheorembg,
	frame hidden,
	boxrule = 0sp,
	borderline west = {2pt}{0pt}{mytheoremfr},
	sharp corners,
	detach title,
	before upper = \tcbtitle\par\smallskip,
	coltitle = mytheoremfr,
	fonttitle = \bfseries\sffamily,
	description font = \mdseries,
	separator sign none,
	segmentation style={solid, mytheoremfr},
}
{th}


\tcbuselibrary{theorems,skins,hooks}
\newtcolorbox{Theoremcon}
{%
	enhanced
	,breakable
	,colback = mytheorembg
	,frame hidden
	,boxrule = 0sp
	,borderline west = {2pt}{0pt}{mytheoremfr}
	,sharp corners
	,description font = \mdseries
	,separator sign none
}

%================================
% Corollery
%================================
\tcbuselibrary{theorems,skins,hooks}
\newtcbtheorem[number within=section]{Corollary}{Corollary}
{%
	enhanced
	,breakable
	,colback = myp!10
	,frame hidden
	,boxrule = 0sp
	,borderline west = {2pt}{0pt}{myp!85!black}
	,sharp corners
	,detach title
	,before upper = \tcbtitle\par\smallskip
	,coltitle = myp!85!black
	,fonttitle = \bfseries\sffamily
	,description font = \mdseries
	,separator sign none
	,segmentation style={solid, myp!85!black}
}
{th}
\tcbuselibrary{theorems,skins,hooks}
\newtcbtheorem[number within=chapter]{corollary}{Corollary}
{%
	enhanced
	,breakable
	,colback = myp!10
	,frame hidden
	,boxrule = 0sp
	,borderline west = {2pt}{0pt}{myp!85!black}
	,sharp corners
	,detach title
	,before upper = \tcbtitle\par\smallskip
	,coltitle = myp!85!black
	,fonttitle = \bfseries\sffamily
	,description font = \mdseries
	,separator sign none
	,segmentation style={solid, myp!85!black}
}
{th}


%================================
% LENMA
%================================

\tcbuselibrary{theorems,skins,hooks}
\newtcbtheorem[number within=section]{Lenma}{Lenma}
{%
	enhanced,
	breakable,
	colback = mylenmabg,
	frame hidden,
	boxrule = 0sp,
	borderline west = {2pt}{0pt}{mylenmafr},
	sharp corners,
	detach title,
	before upper = \tcbtitle\par\smallskip,
	coltitle = mylenmafr,
	fonttitle = \bfseries\sffamily,
	description font = \mdseries,
	separator sign none,
	segmentation style={solid, mylenmafr},
}
{th}

\tcbuselibrary{theorems,skins,hooks}
\newtcbtheorem[number within=chapter]{lenma}{Lenma}
{%
	enhanced,
	breakable,
	colback = mylenmabg,
	frame hidden,
	boxrule = 0sp,
	borderline west = {2pt}{0pt}{mylenmafr},
	sharp corners,
	detach title,
	before upper = \tcbtitle\par\smallskip,
	coltitle = mylenmafr,
	fonttitle = \bfseries\sffamily,
	description font = \mdseries,
	separator sign none,
	segmentation style={solid, mylenmafr},
}
{th}


%================================
% PROPOSITION
%================================

\tcbuselibrary{theorems,skins,hooks}
\newtcbtheorem[number within=section]{Prop}{Proposition}
{%
	enhanced,
	breakable,
	colback = mypropbg,
	frame hidden,
	boxrule = 0sp,
	borderline west = {2pt}{0pt}{mypropfr},
	sharp corners,
	detach title,
	before upper = \tcbtitle\par\smallskip,
	coltitle = mypropfr,
	fonttitle = \bfseries\sffamily,
	description font = \mdseries,
	separator sign none,
	segmentation style={solid, mypropfr},
}
{th}

\tcbuselibrary{theorems,skins,hooks}
\newtcbtheorem[number within=chapter]{prop}{Proposition}
{%
	enhanced,
	breakable,
	colback = mypropbg,
	frame hidden,
	boxrule = 0sp,
	borderline west = {2pt}{0pt}{mypropfr},
	sharp corners,
	detach title,
	before upper = \tcbtitle\par\smallskip,
	coltitle = mypropfr,
	fonttitle = \bfseries\sffamily,
	description font = \mdseries,
	separator sign none,
	segmentation style={solid, mypropfr},
}
{th}


%================================
% CLAIM
%================================

\tcbuselibrary{theorems,skins,hooks}
\newtcbtheorem[number within=section]{claim}{Claim}
{%
	enhanced
	,breakable
	,colback = myg!10
	,frame hidden
	,boxrule = 0sp
	,borderline west = {2pt}{0pt}{myg}
	,sharp corners
	,detach title
	,before upper = \tcbtitle\par\smallskip
	,coltitle = myg!85!black
	,fonttitle = \bfseries\sffamily
	,description font = \mdseries
	,separator sign none
	,segmentation style={solid, myg!85!black}
}
{th}



%================================
% Exercise
%================================

\tcbuselibrary{theorems,skins,hooks}
\newtcbtheorem[number within=section]{Exercise}{Exercise}
{%
	enhanced,
	breakable,
	colback = myexercisebg,
	frame hidden,
	boxrule = 0sp,
	borderline west = {2pt}{0pt}{myexercisefg},
	sharp corners,
	detach title,
	before upper = \tcbtitle\par\smallskip,
	coltitle = myexercisefg,
	fonttitle = \bfseries\sffamily,
	description font = \mdseries,
	separator sign none,
	segmentation style={solid, myexercisefg},
}
{th}

\tcbuselibrary{theorems,skins,hooks}
\newtcbtheorem[number within=chapter]{exercise}{Exercise}
{%
	enhanced,
	breakable,
	colback = myexercisebg,
	frame hidden,
	boxrule = 0sp,
	borderline west = {2pt}{0pt}{myexercisefg},
	sharp corners,
	detach title,
	before upper = \tcbtitle\par\smallskip,
	coltitle = myexercisefg,
	fonttitle = \bfseries\sffamily,
	description font = \mdseries,
	separator sign none,
	segmentation style={solid, myexercisefg},
}
{th}

%================================
% EXAMPLE BOX
%================================

\newtcbtheorem[number within=section]{Example}{Example}
{%
	colback = myexamplebg
	,breakable
	,colframe = myexamplefr
	,coltitle = myexampleti
	,boxrule = 1pt
	,sharp corners
	,detach title
	,before upper=\tcbtitle\par\smallskip
	,fonttitle = \bfseries
	,description font = \mdseries
	,separator sign none
	,description delimiters parenthesis
}
{ex}

\newtcbtheorem[number within=chapter]{example}{Example}
{%
	colback = myexamplebg
	,breakable
	,colframe = myexamplefr
	,coltitle = myexampleti
	,boxrule = 1pt
	,sharp corners
	,detach title
	,before upper=\tcbtitle\par\smallskip
	,fonttitle = \bfseries
	,description font = \mdseries
	,separator sign none
	,description delimiters parenthesis
}
{ex}

%================================
% DEFINITION BOX
%================================

\newtcbtheorem[number within=section]{Definition}{Definition}{enhanced,
	before skip=2mm,after skip=2mm, colback=red!5,colframe=red!80!black,boxrule=0.5mm,
	attach boxed title to top left={xshift=1cm,yshift*=1mm-\tcboxedtitleheight}, varwidth boxed title*=-3cm,
	boxed title style={frame code={
					\path[fill=tcbcolback]
					([yshift=-1mm,xshift=-1mm]frame.north west)
					arc[start angle=0,end angle=180,radius=1mm]
					([yshift=-1mm,xshift=1mm]frame.north east)
					arc[start angle=180,end angle=0,radius=1mm];
					\path[left color=tcbcolback!60!black,right color=tcbcolback!60!black,
						middle color=tcbcolback!80!black]
					([xshift=-2mm]frame.north west) -- ([xshift=2mm]frame.north east)
					[rounded corners=1mm]-- ([xshift=1mm,yshift=-1mm]frame.north east)
					-- (frame.south east) -- (frame.south west)
					-- ([xshift=-1mm,yshift=-1mm]frame.north west)
					[sharp corners]-- cycle;
				},interior engine=empty,
		},
	fonttitle=\bfseries,
	title={#2},#1}{def}
\newtcbtheorem[number within=chapter]{definition}{Definition}{enhanced,
	before skip=2mm,after skip=2mm, colback=red!5,colframe=red!80!black,boxrule=0.5mm,
	attach boxed title to top left={xshift=1cm,yshift*=1mm-\tcboxedtitleheight}, varwidth boxed title*=-3cm,
	boxed title style={frame code={
					\path[fill=tcbcolback]
					([yshift=-1mm,xshift=-1mm]frame.north west)
					arc[start angle=0,end angle=180,radius=1mm]
					([yshift=-1mm,xshift=1mm]frame.north east)
					arc[start angle=180,end angle=0,radius=1mm];
					\path[left color=tcbcolback!60!black,right color=tcbcolback!60!black,
						middle color=tcbcolback!80!black]
					([xshift=-2mm]frame.north west) -- ([xshift=2mm]frame.north east)
					[rounded corners=1mm]-- ([xshift=1mm,yshift=-1mm]frame.north east)
					-- (frame.south east) -- (frame.south west)
					-- ([xshift=-1mm,yshift=-1mm]frame.north west)
					[sharp corners]-- cycle;
				},interior engine=empty,
		},
	fonttitle=\bfseries,
	title={#2},#1}{def}



%================================
% Solution BOX
%================================

\makeatletter
\newtcbtheorem{question}{Question}{enhanced,
	breakable,
	colback=white,
	colframe=myb!80!black,
	attach boxed title to top left={yshift*=-\tcboxedtitleheight},
	fonttitle=\bfseries,
	title={#2},
	boxed title size=title,
	boxed title style={%
			sharp corners,
			rounded corners=northwest,
			colback=tcbcolframe,
			boxrule=0pt,
		},
	underlay boxed title={%
			\path[fill=tcbcolframe] (title.south west)--(title.south east)
			to[out=0, in=180] ([xshift=5mm]title.east)--
			(title.center-|frame.east)
			[rounded corners=\kvtcb@arc] |-
			(frame.north) -| cycle;
		},
	#1
}{def}
\makeatother

%================================
% SOLUTION BOX
%================================

\makeatletter
\newtcolorbox{solution}{enhanced,
	breakable,
	colback=white,
	colframe=myg!80!black,
	attach boxed title to top left={yshift*=-\tcboxedtitleheight},
	title=Solution,
	boxed title size=title,
	boxed title style={%
			sharp corners,
			rounded corners=northwest,
			colback=tcbcolframe,
			boxrule=0pt,
		},
	underlay boxed title={%
			\path[fill=tcbcolframe] (title.south west)--(title.south east)
			to[out=0, in=180] ([xshift=5mm]title.east)--
			(title.center-|frame.east)
			[rounded corners=\kvtcb@arc] |-
			(frame.north) -| cycle;
		},
}
\makeatother

%================================
% Question BOX
%================================

\makeatletter
\newtcbtheorem{qstion}{Question}{enhanced,
	breakable,
	colback=white,
	colframe=mygr,
	attach boxed title to top left={yshift*=-\tcboxedtitleheight},
	fonttitle=\bfseries,
	title={#2},
	boxed title size=title,
	boxed title style={%
			sharp corners,
			rounded corners=northwest,
			colback=tcbcolframe,
			boxrule=0pt,
		},
	underlay boxed title={%
			\path[fill=tcbcolframe] (title.south west)--(title.south east)
			to[out=0, in=180] ([xshift=5mm]title.east)--
			(title.center-|frame.east)
			[rounded corners=\kvtcb@arc] |-
			(frame.north) -| cycle;
		},
	#1
}{def}
\makeatother

\newtcbtheorem[number within=chapter]{wconc}{Wrong Concept}{
	breakable,
	enhanced,
	colback=white,
	colframe=myr,
	arc=0pt,
	outer arc=0pt,
	fonttitle=\bfseries\sffamily\large,
	colbacktitle=myr,
	attach boxed title to top left={},
	boxed title style={
			enhanced,
			skin=enhancedfirst jigsaw,
			arc=3pt,
			bottom=0pt,
			interior style={fill=myr}
		},
	#1
}{def}



%================================
% NOTE BOX
%================================

\usetikzlibrary{arrows,calc,shadows.blur}
\tcbuselibrary{skins}
\newtcolorbox{note}[1][]{%
	enhanced jigsaw,
	colback=gray!20!white,%
	colframe=gray!80!black,
	size=small,
	boxrule=1pt,
	title=\textbf{Note:},
	halign title=flush center,
	coltitle=black,
	breakable,
	drop shadow=black!50!white,
	attach boxed title to top left={xshift=1cm,yshift=-\tcboxedtitleheight/2,yshifttext=-\tcboxedtitleheight/2},
	minipage boxed title=1.5cm,
	boxed title style={%
			colback=white,
			size=fbox,
			boxrule=1pt,
			boxsep=2pt,
			underlay={%
					\coordinate (dotA) at ($(interior.west) + (-0.5pt,0)$);
					\coordinate (dotB) at ($(interior.east) + (0.5pt,0)$);
					\begin{scope}
						\clip (interior.north west) rectangle ([xshift=3ex]interior.east);
						\filldraw [white, blur shadow={shadow opacity=60, shadow yshift=-.75ex}, rounded corners=2pt] (interior.north west) rectangle (interior.south east);
					\end{scope}
					\begin{scope}[gray!80!black]
						\fill (dotA) circle (2pt);
						\fill (dotB) circle (2pt);
					\end{scope}
				},
		},
	#1,
}

%%%%%%%%%%%%%%%%%%%%%%%%%%%%%%
% SELF MADE COMMANDS
%%%%%%%%%%%%%%%%%%%%%%%%%%%%%%


\newcommand{\thm}[2]{\begin{Theorem}{#1}{}#2\end{Theorem}}
\newcommand{\cor}[2]{\begin{Corollary}{#1}{}#2\end{Corollary}}
\newcommand{\mlenma}[2]{\begin{Lenma}{#1}{}#2\end{Lenma}}
\newcommand{\mprop}[2]{\begin{Prop}{#1}{}#2\end{Prop}}
\newcommand{\clm}[3]{\begin{claim}{#1}{#2}#3\end{claim}}
\newcommand{\wc}[2]{\begin{wconc}{#1}{}\setlength{\parindent}{1cm}#2\end{wconc}}
\newcommand{\thmcon}[1]{\begin{Theoremcon}{#1}\end{Theoremcon}}
\newcommand{\ex}[2]{\begin{Example}{#1}{}#2\end{Example}}
\newcommand{\dfn}[2]{\begin{Definition}[colbacktitle=red!75!black]{#1}{}#2\end{Definition}}
\newcommand{\dfnc}[2]{\begin{definition}[colbacktitle=red!75!black]{#1}{}#2\end{definition}}
\newcommand{\qs}[2]{\begin{question}{#1}{}#2\end{question}}
\newcommand{\pf}[2]{\begin{myproof}[#1]#2\end{myproof}}
\newcommand{\nt}[1]{\begin{note}#1\end{note}}

\newcommand*\circled[1]{\tikz[baseline=(char.base)]{
		\node[shape=circle,draw,inner sep=1pt] (char) {#1};}}
\newcommand\getcurrentref[1]{%
	\ifnumequal{\value{#1}}{0}
	{??}
	{\the\value{#1}}%
}
\newcommand{\getCurrentSectionNumber}{\getcurrentref{section}}
\newenvironment{myproof}[1][\proofname]{%
	\proof[\bfseries #1: ]%
}{\endproof}

\newcommand{\mclm}[2]{\begin{myclaim}[#1]#2\end{myclaim}}
\newenvironment{myclaim}[1][\claimname]{\proof[\bfseries #1: ]}{}

\newcounter{mylabelcounter}

\makeatletter
\newcommand{\setword}[2]{%
	\phantomsection
	#1\def\@currentlabel{\unexpanded{#1}}\label{#2}%
}
\makeatother




\tikzset{
	symbol/.style={
			draw=none,
			every to/.append style={
					edge node={node [sloped, allow upside down, auto=false]{$#1$}}}
		}
}


% deliminators
\DeclarePairedDelimiter{\abs}{\lvert}{\rvert}
\DeclarePairedDelimiter{\norm}{\lVert}{\rVert}

\DeclarePairedDelimiter{\ceil}{\lceil}{\rceil}
\DeclarePairedDelimiter{\floor}{\lfloor}{\rfloor}
\DeclarePairedDelimiter{\round}{\lfloor}{\rceil}

\newsavebox\diffdbox
\newcommand{\slantedromand}{{\mathpalette\makesl{d}}}
\newcommand{\makesl}[2]{%
\begingroup
\sbox{\diffdbox}{$\mathsurround=0pt#1\mathrm{#2}$}%
\pdfsave
\pdfsetmatrix{1 0 0.2 1}%
\rlap{\usebox{\diffdbox}}%
\pdfrestore
\hskip\wd\diffdbox
\endgroup
}
\newcommand{\dd}[1][]{\ensuremath{\mathop{}\!\ifstrempty{#1}{%
\slantedromand\@ifnextchar^{\hspace{0.2ex}}{\hspace{0.1ex}}}%
{\slantedromand\hspace{0.2ex}^{#1}}}}
\ProvideDocumentCommand\dv{o m g}{%
  \ensuremath{%
    \IfValueTF{#3}{%
      \IfNoValueTF{#1}{%
        \frac{\dd #2}{\dd #3}%
      }{%
        \frac{\dd^{#1} #2}{\dd #3^{#1}}%
      }%
    }{%
      \IfNoValueTF{#1}{%
        \frac{\dd}{\dd #2}%
      }{%
        \frac{\dd^{#1}}{\dd #2^{#1}}%
      }%
    }%
  }%
}
\providecommand*{\pdv}[3][]{\frac{\partial^{#1}#2}{\partial#3^{#1}}}
%  - others
\DeclareMathOperator{\Lap}{\mathcal{L}}
\DeclareMathOperator{\Var}{Var} % varience
\DeclareMathOperator{\Cov}{Cov} % covarience
\DeclareMathOperator{\E}{E} % expected

% Since the amsthm package isn't loaded

% I prefer the slanted \leq
\let\oldleq\leq % save them in case they're every wanted
\let\oldgeq\geq
\renewcommand{\leq}{\leqslant}
\renewcommand{\geq}{\geqslant}

% % redefine matrix env to allow for alignment, use r as default
% \renewcommand*\env@matrix[1][r]{\hskip -\arraycolsep
%     \let\@ifnextchar\new@ifnextchar
%     \array{*\c@MaxMatrixCols #1}}


%\usepackage{framed}
%\usepackage{titletoc}
%\usepackage{etoolbox}
%\usepackage{lmodern}


%\patchcmd{\tableofcontents}{\contentsname}{\sffamily\contentsname}{}{}

%\renewenvironment{leftbar}
%{\def\FrameCommand{\hspace{6em}%
%		{\color{myyellow}\vrule width 2pt depth 6pt}\hspace{1em}}%
%	\MakeFramed{\parshape 1 0cm \dimexpr\textwidth-6em\relax\FrameRestore}\vskip2pt%
%}
%{\endMakeFramed}

%\titlecontents{chapter}
%[0em]{\vspace*{2\baselineskip}}
%{\parbox{4.5em}{%
%		\hfill\Huge\sffamily\bfseries\color{myred}\thecontentspage}%
%	\vspace*{-2.3\baselineskip}\leftbar\textsc{\small\chaptername~\thecontentslabel}\\\sffamily}
%{}{\endleftbar}
%\titlecontents{section}
%[8.4em]
%{\sffamily\contentslabel{3em}}{}{}
%{\hspace{0.5em}\nobreak\itshape\color{myred}\contentspage}
%\titlecontents{subsection}
%[8.4em]
%{\sffamily\contentslabel{3em}}{}{}  
%{\hspace{0.5em}\nobreak\itshape\color{myred}\contentspage}



%%%%%%%%%%%%%%%%%%%%%%%%%%%%%%%%%%%%%%%%%%%
% TABLE OF CONTENTS
%%%%%%%%%%%%%%%%%%%%%%%%%%%%%%%%%%%%%%%%%%%

\usepackage{tikz}
\definecolor{doc}{RGB}{0,60,110}
\usepackage{titletoc}
\contentsmargin{0cm}
\titlecontents{chapter}[3.7pc]
{\addvspace{30pt}%
	\begin{tikzpicture}[remember picture, overlay]%
		\draw[fill=doc!60,draw=doc!60] (-7,-.1) rectangle (-0.9,.5);%
		\pgftext[left,x=-3.5cm,y=0.2cm]{\color{white}\Large\sc\bfseries Chapter\ \thecontentslabel};%
	\end{tikzpicture}\color{doc!60}\large\sc\bfseries}%
{}
{}
{\;\titlerule\;\large\sc\bfseries Page \thecontentspage
	\begin{tikzpicture}[remember picture, overlay]
		\draw[fill=doc!60,draw=doc!60] (2pt,0) rectangle (4,0.1pt);
	\end{tikzpicture}}%
\titlecontents{section}[3.7pc]
{\addvspace{2pt}}
{\contentslabel[\thecontentslabel]{2pc}}
{}
{\hfill\small \thecontentspage}
[]
\titlecontents*{subsection}[3.7pc]
{\addvspace{-1pt}\small}
{}
{}
{\ --- \small\thecontentspage}
[ \textbullet\ ][]

\makeatletter
\renewcommand{\tableofcontents}{%
	\chapter*{%
	  \vspace*{-20\p@}%
	  \begin{tikzpicture}[remember picture, overlay]%
		  \pgftext[right,x=15cm,y=0.2cm]{\color{doc!60}\Huge\sc\bfseries \contentsname};%
		  \draw[fill=doc!60,draw=doc!60] (13,-.75) rectangle (20,1);%
		  \clip (13,-.75) rectangle (20,1);
		  \pgftext[right,x=15cm,y=0.2cm]{\color{white}\Huge\sc\bfseries \contentsname};%
	  \end{tikzpicture}}%
	\@starttoc{toc}}
\makeatother


% Comandi di Bastiality %
\newcommand{\innerproduct}[2]{\langle #1, #2 \rangle}
\newcommand{\generators}[1]{\langle #1 \rangle}
\newcommand{\powerset}[0]{\mathcal{P}}
\newcommand{\indip}{\perp \!\!\! \perp}

% Comandi di GioLaPalma %
\newcommand{\der}[0]{\implies^*}
\newcommand{\red}[1]{\textcolor{red}{#1}}
\newcommand{\esempio}[1]{\ex{}{#1}}
\newcommand{\teorema}[1]{\thm{}{#1}}
\newcommand{\dimostrazione}[1]{\pf{Dimostrazione}{#1}}
\setlength{\parindent}{0pt}
\newcommand{\osservazione}[1]{\nt{#1}}
\newcommand{\inside}[1]{
  \langle #1 \rangle
}
      

%From M275 "Topology" at SJSU
\newcommand{\id}{\mathrm{id}}
\newcommand{\taking}[1]{\xrightarrow{#1}}
\newcommand{\inv}{^{-1}}

%From M170 "Introduction to Graph Theory" at SJSU
\DeclareMathOperator{\diam}{diam}
\DeclareMathOperator{\ord}{ord}
\newcommand{\defeq}{\overset{\mathrm{def}}{=}}

%From the USAMO .tex files
\newcommand{\ts}{\textsuperscript}
\newcommand{\dg}{^\circ}
\newcommand{\ii}{\item}

% % From Math 55 and Math 145 at Harvard
% \newenvironment{subproof}[1][Proof]{%
% \begin{proof}[#1] \renewcommand{\qedsymbol}{$\blacksquare$}}%
% {\end{proof}}

\newcommand{\liff}{\leftrightarrow}
\newcommand{\lthen}{\rightarrow}
\newcommand{\opname}{\operatorname}
\newcommand{\surjto}{\twoheadrightarrow}
\newcommand{\injto}{\hookrightarrow}
\newcommand{\On}{\mathrm{On}} % ordinals
\DeclareMathOperator{\img}{im} % Image
\DeclareMathOperator{\Img}{Im} % Image
\DeclareMathOperator{\coker}{coker} % Cokernel
\DeclareMathOperator{\Coker}{Coker} % Cokernel
\DeclareMathOperator{\Ker}{Ker} % Kernel
\DeclareMathOperator{\rank}{rank}
\DeclareMathOperator{\Spec}{Spec} % spectrum
\DeclareMathOperator{\Tr}{Tr} % trace
\DeclareMathOperator{\pr}{pr} % projection
\DeclareMathOperator{\ext}{ext} % extension
\DeclareMathOperator{\pred}{pred} % predecessor
\DeclareMathOperator{\dom}{dom} % domain
\DeclareMathOperator{\ran}{ran} % range
\DeclareMathOperator{\Hom}{Hom} % homomorphism
\DeclareMathOperator{\Mor}{Mor} % morphisms
\DeclareMathOperator{\End}{End} % endomorphism

\newcommand{\eps}{\epsilon}
\newcommand{\veps}{\varepsilon}
\newcommand{\ol}{\overline}
\newcommand{\ul}{\underline}
\newcommand{\wt}{\widetilde}
\newcommand{\wh}{\widehat}
\newcommand{\vocab}[1]{\textbf{\color{blue} #1}}
\providecommand{\half}{\frac{1}{2}}
\newcommand{\dang}{\measuredangle} %% Directed angle
\newcommand{\ray}[1]{\overrightarrow{#1}}
\newcommand{\seg}[1]{\overline{#1}}
\newcommand{\arc}[1]{\wideparen{#1}}
\DeclareMathOperator{\cis}{cis}
\DeclareMathOperator*{\lcm}{lcm}
\DeclareMathOperator*{\argmin}{arg min}
\DeclareMathOperator*{\argmax}{arg max}
\newcommand{\cycsum}{\sum_{\mathrm{cyc}}}
\newcommand{\symsum}{\sum_{\mathrm{sym}}}
\newcommand{\cycprod}{\prod_{\mathrm{cyc}}}
\newcommand{\symprod}{\prod_{\mathrm{sym}}}
\newcommand{\Qed}{\begin{flushright}\qed\end{flushright}}
\newcommand{\parinn}{\setlength{\parindent}{1cm}}
\newcommand{\parinf}{\setlength{\parindent}{0cm}}
% \newcommand{\norm}{\|\cdot\|}
\newcommand{\inorm}{\norm_{\infty}}
\newcommand{\opensets}{\{V_{\alpha}\}_{\alpha\in I}}
\newcommand{\oset}{V_{\alpha}}
\newcommand{\opset}[1]{V_{\alpha_{#1}}}
\newcommand{\lub}{\text{lub}}
\newcommand{\del}[2]{\frac{\partial #1}{\partial #2}}
\newcommand{\Del}[3]{\frac{\partial^{#1} #2}{\partial^{#1} #3}}
\newcommand{\deld}[2]{\dfrac{\partial #1}{\partial #2}}
\newcommand{\Deld}[3]{\dfrac{\partial^{#1} #2}{\partial^{#1} #3}}
\newcommand{\lm}{\lambda}
\newcommand{\uin}{\mathbin{\rotatebox[origin=c]{90}{$\in$}}}
\newcommand{\usubset}{\mathbin{\rotatebox[origin=c]{90}{$\subset$}}}
\newcommand{\lt}{\left}
\newcommand{\rt}{\right}
\newcommand{\bs}[1]{\boldsymbol{#1}}
\newcommand{\exs}{\exists}
\newcommand{\st}{\strut}
\newcommand{\dps}[1]{\displaystyle{#1}}

\newcommand{\sol}{\setlength{\parindent}{0cm}\textbf{\textit{Solution:}}\setlength{\parindent}{1cm} }
\newcommand{\solve}[1]{\setlength{\parindent}{0cm}\textbf{\textit{Solution: }}\setlength{\parindent}{1cm}#1 \Qed}

% Things Lie
\newcommand{\kb}{\mathfrak b}
\newcommand{\kg}{\mathfrak g}
\newcommand{\kh}{\mathfrak h}
\newcommand{\kn}{\mathfrak n}
\newcommand{\ku}{\mathfrak u}
\newcommand{\kz}{\mathfrak z}
\DeclareMathOperator{\Ext}{Ext} % Ext functor
\DeclareMathOperator{\Tor}{Tor} % Tor functor
\newcommand{\gl}{\opname{\mathfrak{gl}}} % frak gl group
\renewcommand{\sl}{\opname{\mathfrak{sl}}} % frak sl group chktex 6

% More script letters etc.
\newcommand{\SA}{\mathcal A}
\newcommand{\SB}{\mathcal B}
\newcommand{\SC}{\mathcal C}
\newcommand{\SF}{\mathcal F}
\newcommand{\SG}{\mathcal G}
\newcommand{\SH}{\mathcal H}
\newcommand{\OO}{\mathcal O}

\newcommand{\SCA}{\mathscr A}
\newcommand{\SCB}{\mathscr B}
\newcommand{\SCC}{\mathscr C}
\newcommand{\SCD}{\mathscr D}
\newcommand{\SCE}{\mathscr E}
\newcommand{\SCF}{\mathscr F}
\newcommand{\SCG}{\mathscr G}
\newcommand{\SCH}{\mathscr H}

% Mathfrak primes
\newcommand{\km}{\mathfrak m}
\newcommand{\kp}{\mathfrak p}
\newcommand{\kq}{\mathfrak q}

% number sets
\newcommand{\RR}[1][]{\ensuremath{\ifstrempty{#1}{\mathbb{R}}{\mathbb{R}^{#1}}}}
\newcommand{\NN}[1][]{\ensuremath{\ifstrempty{#1}{\mathbb{N}}{\mathbb{N}^{#1}}}}
\newcommand{\ZZ}[1][]{\ensuremath{\ifstrempty{#1}{\mathbb{Z}}{\mathbb{Z}^{#1}}}}
\newcommand{\QQ}[1][]{\ensuremath{\ifstrempty{#1}{\mathbb{Q}}{\mathbb{Q}^{#1}}}}
\newcommand{\CC}[1][]{\ensuremath{\ifstrempty{#1}{\mathbb{C}}{\mathbb{C}^{#1}}}}
\newcommand{\PP}[1][]{\ensuremath{\ifstrempty{#1}{\mathbb{P}}{\mathbb{P}^{#1}}}}
\newcommand{\HH}[1][]{\ensuremath{\ifstrempty{#1}{\mathbb{H}}{\mathbb{H}^{#1}}}}
\newcommand{\FF}[1][]{\ensuremath{\ifstrempty{#1}{\mathbb{F}}{\mathbb{F}^{#1}}}}
% expected value
\newcommand{\EE}{\ensuremath{\mathbb{E}}}
\newcommand{\charin}{\text{ char }}
\DeclareMathOperator{\sign}{sign}
\DeclareMathOperator{\Aut}{Aut}
\DeclareMathOperator{\Inn}{Inn}
\DeclareMathOperator{\Syl}{Syl}
\DeclareMathOperator{\Gal}{Gal}
\DeclareMathOperator{\GL}{GL} % General linear group
\DeclareMathOperator{\SL}{SL} % Special linear group

%---------------------------------------
% BlackBoard Math Fonts :-
%---------------------------------------

%Captital Letters
\newcommand{\bbA}{\mathbb{A}}	\newcommand{\bbB}{\mathbb{B}}
\newcommand{\bbC}{\mathbb{C}}	\newcommand{\bbD}{\mathbb{D}}
\newcommand{\bbE}{\mathbb{E}}	\newcommand{\bbF}{\mathbb{F}}
\newcommand{\bbG}{\mathbb{G}}	\newcommand{\bbH}{\mathbb{H}}
\newcommand{\bbI}{\mathbb{I}}	\newcommand{\bbJ}{\mathbb{J}}
\newcommand{\bbK}{\mathbb{K}}	\newcommand{\bbL}{\mathbb{L}}
\newcommand{\bbM}{\mathbb{M}}	\newcommand{\bbN}{\mathbb{N}}
\newcommand{\bbO}{\mathbb{O}}	\newcommand{\bbP}{\mathbb{P}}
\newcommand{\bbQ}{\mathbb{Q}}	\newcommand{\bbR}{\mathbb{R}}
\newcommand{\bbS}{\mathbb{S}}	\newcommand{\bbT}{\mathbb{T}}
\newcommand{\bbU}{\mathbb{U}}	\newcommand{\bbV}{\mathbb{V}}
\newcommand{\bbW}{\mathbb{W}}	\newcommand{\bbX}{\mathbb{X}}
\newcommand{\bbY}{\mathbb{Y}}	\newcommand{\bbZ}{\mathbb{Z}}

%---------------------------------------
% MathCal Fonts :-
%---------------------------------------

%Captital Letters
\newcommand{\mcA}{\mathcal{A}}	\newcommand{\mcB}{\mathcal{B}}
\newcommand{\mcC}{\mathcal{C}}	\newcommand{\mcD}{\mathcal{D}}
\newcommand{\mcE}{\mathcal{E}}	\newcommand{\mcF}{\mathcal{F}}
\newcommand{\mcG}{\mathcal{G}}	\newcommand{\mcH}{\mathcal{H}}
\newcommand{\mcI}{\mathcal{I}}	\newcommand{\mcJ}{\mathcal{J}}
\newcommand{\mcK}{\mathcal{K}}	\newcommand{\mcL}{\mathcal{L}}
\newcommand{\mcM}{\mathcal{M}}	\newcommand{\mcN}{\mathcal{N}}
\newcommand{\mcO}{\mathcal{O}}	\newcommand{\mcP}{\mathcal{P}}
\newcommand{\mcQ}{\mathcal{Q}}	\newcommand{\mcR}{\mathcal{R}}
\newcommand{\mcS}{\mathcal{S}}	\newcommand{\mcT}{\mathcal{T}}
\newcommand{\mcU}{\mathcal{U}}	\newcommand{\mcV}{\mathcal{V}}
\newcommand{\mcW}{\mathcal{W}}	\newcommand{\mcX}{\mathcal{X}}
\newcommand{\mcY}{\mathcal{Y}}	\newcommand{\mcZ}{\mathcal{Z}}


%---------------------------------------
% Bold Math Fonts :-
%---------------------------------------

%Captital Letters
\newcommand{\bmA}{\boldsymbol{A}}	\newcommand{\bmB}{\boldsymbol{B}}
\newcommand{\bmC}{\boldsymbol{C}}	\newcommand{\bmD}{\boldsymbol{D}}
\newcommand{\bmE}{\boldsymbol{E}}	\newcommand{\bmF}{\boldsymbol{F}}
\newcommand{\bmG}{\boldsymbol{G}}	\newcommand{\bmH}{\boldsymbol{H}}
\newcommand{\bmI}{\boldsymbol{I}}	\newcommand{\bmJ}{\boldsymbol{J}}
\newcommand{\bmK}{\boldsymbol{K}}	\newcommand{\bmL}{\boldsymbol{L}}
\newcommand{\bmM}{\boldsymbol{M}}	\newcommand{\bmN}{\boldsymbol{N}}
\newcommand{\bmO}{\boldsymbol{O}}	\newcommand{\bmP}{\boldsymbol{P}}
\newcommand{\bmQ}{\boldsymbol{Q}}	\newcommand{\bmR}{\boldsymbol{R}}
\newcommand{\bmS}{\boldsymbol{S}}	\newcommand{\bmT}{\boldsymbol{T}}
\newcommand{\bmU}{\boldsymbol{U}}	\newcommand{\bmV}{\boldsymbol{V}}
\newcommand{\bmW}{\boldsymbol{W}}	\newcommand{\bmX}{\boldsymbol{X}}
\newcommand{\bmY}{\boldsymbol{Y}}	\newcommand{\bmZ}{\boldsymbol{Z}}
%Small Letters
\newcommand{\bma}{\boldsymbol{a}}	\newcommand{\bmb}{\boldsymbol{b}}
\newcommand{\bmc}{\boldsymbol{c}}	\newcommand{\bmd}{\boldsymbol{d}}
\newcommand{\bme}{\boldsymbol{e}}	\newcommand{\bmf}{\boldsymbol{f}}
\newcommand{\bmg}{\boldsymbol{g}}	\newcommand{\bmh}{\boldsymbol{h}}
\newcommand{\bmi}{\boldsymbol{i}}	\newcommand{\bmj}{\boldsymbol{j}}
\newcommand{\bmk}{\boldsymbol{k}}	\newcommand{\bml}{\boldsymbol{l}}
\newcommand{\bmm}{\boldsymbol{m}}	\newcommand{\bmn}{\boldsymbol{n}}
\newcommand{\bmo}{\boldsymbol{o}}	\newcommand{\bmp}{\boldsymbol{p}}
\newcommand{\bmq}{\boldsymbol{q}}	\newcommand{\bmr}{\boldsymbol{r}}
\newcommand{\bms}{\boldsymbol{s}}	\newcommand{\bmt}{\boldsymbol{t}}
\newcommand{\bmu}{\boldsymbol{u}}	\newcommand{\bmv}{\boldsymbol{v}}
\newcommand{\bmw}{\boldsymbol{w}}	\newcommand{\bmx}{\boldsymbol{x}}
\newcommand{\bmy}{\boldsymbol{y}}	\newcommand{\bmz}{\boldsymbol{z}}

%---------------------------------------
% Scr Math Fonts :-
%---------------------------------------

\newcommand{\sA}{{\mathscr{A}}}   \newcommand{\sB}{{\mathscr{B}}}
\newcommand{\sC}{{\mathscr{C}}}   \newcommand{\sD}{{\mathscr{D}}}
\newcommand{\sE}{{\mathscr{E}}}   \newcommand{\sF}{{\mathscr{F}}}
\newcommand{\sG}{{\mathscr{G}}}   \newcommand{\sH}{{\mathscr{H}}}
\newcommand{\sI}{{\mathscr{I}}}   \newcommand{\sJ}{{\mathscr{J}}}
\newcommand{\sK}{{\mathscr{K}}}   \newcommand{\sL}{{\mathscr{L}}}
\newcommand{\sM}{{\mathscr{M}}}   \newcommand{\sN}{{\mathscr{N}}}
\newcommand{\sO}{{\mathscr{O}}}   \newcommand{\sP}{{\mathscr{P}}}
\newcommand{\sQ}{{\mathscr{Q}}}   \newcommand{\sR}{{\mathscr{R}}}
\newcommand{\sS}{{\mathscr{S}}}   \newcommand{\sT}{{\mathscr{T}}}
\newcommand{\sU}{{\mathscr{U}}}   \newcommand{\sV}{{\mathscr{V}}}
\newcommand{\sW}{{\mathscr{W}}}   \newcommand{\sX}{{\mathscr{X}}}
\newcommand{\sY}{{\mathscr{Y}}}   \newcommand{\sZ}{{\mathscr{Z}}}


%---------------------------------------
% Math Fraktur Font
%---------------------------------------

%Captital Letters
\newcommand{\mfA}{\mathfrak{A}}	\newcommand{\mfB}{\mathfrak{B}}
\newcommand{\mfC}{\mathfrak{C}}	\newcommand{\mfD}{\mathfrak{D}}
\newcommand{\mfE}{\mathfrak{E}}	\newcommand{\mfF}{\mathfrak{F}}
\newcommand{\mfG}{\mathfrak{G}}	\newcommand{\mfH}{\mathfrak{H}}
\newcommand{\mfI}{\mathfrak{I}}	\newcommand{\mfJ}{\mathfrak{J}}
\newcommand{\mfK}{\mathfrak{K}}	\newcommand{\mfL}{\mathfrak{L}}
\newcommand{\mfM}{\mathfrak{M}}	\newcommand{\mfN}{\mathfrak{N}}
\newcommand{\mfO}{\mathfrak{O}}	\newcommand{\mfP}{\mathfrak{P}}
\newcommand{\mfQ}{\mathfrak{Q}}	\newcommand{\mfR}{\mathfrak{R}}
\newcommand{\mfS}{\mathfrak{S}}	\newcommand{\mfT}{\mathfrak{T}}
\newcommand{\mfU}{\mathfrak{U}}	\newcommand{\mfV}{\mathfrak{V}}
\newcommand{\mfW}{\mathfrak{W}}	\newcommand{\mfX}{\mathfrak{X}}
\newcommand{\mfY}{\mathfrak{Y}}	\newcommand{\mfZ}{\mathfrak{Z}}
%Small Letters
\newcommand{\mfa}{\mathfrak{a}}	\newcommand{\mfb}{\mathfrak{b}}
\newcommand{\mfc}{\mathfrak{c}}	\newcommand{\mfd}{\mathfrak{d}}
\newcommand{\mfe}{\mathfrak{e}}	\newcommand{\mff}{\mathfrak{f}}
\newcommand{\mfg}{\mathfrak{g}}	\newcommand{\mfh}{\mathfrak{h}}
\newcommand{\mfi}{\mathfrak{i}}	\newcommand{\mfj}{\mathfrak{j}}
\newcommand{\mfk}{\mathfrak{k}}	\newcommand{\mfl}{\mathfrak{l}}
\newcommand{\mfm}{\mathfrak{m}}	\newcommand{\mfn}{\mathfrak{n}}
\newcommand{\mfo}{\mathfrak{o}}	\newcommand{\mfp}{\mathfrak{p}}
\newcommand{\mfq}{\mathfrak{q}}	\newcommand{\mfr}{\mathfrak{r}}
\newcommand{\mfs}{\mathfrak{s}}	\newcommand{\mft}{\mathfrak{t}}
\newcommand{\mfu}{\mathfrak{u}}	\newcommand{\mfv}{\mathfrak{v}}
\newcommand{\mfw}{\mathfrak{w}}	\newcommand{\mfx}{\mathfrak{x}}
\newcommand{\mfy}{\mathfrak{y}}	\newcommand{\mfz}{\mathfrak{z}}


\usepackage[T1]{fontenc}
\usepackage[utf8]{inputenc}
\usepackage{xcolor}
\usepackage{listings}
\usepackage{array}
\usepackage{longtable}
\PassOptionsToPackage{hyphens}{url}
\usepackage{url}
\Urlmuskip=0mu plus 1mu

%==============================================================
\definecolor{codebg}{HTML}{F7F7FB}
\definecolor{codekw}{HTML}{2D6FB1}
\definecolor{codecm}{HTML}{388C46}
\definecolor{codestr}{HTML}{C74440}
\definecolor{coderule}{HTML}{CCCCCC}

\lstset{
    basicstyle=\ttfamily\footnotesize,
    keywordstyle=\color{codekw}\bfseries,
    commentstyle=\color{codecm}\itshape,
    stringstyle=\color{codestr},
    numberstyle=\tiny\color{gray},
    showstringspaces=false,
    breaklines=true,
    breakatwhitespace=false,
    frame=single,
    rulecolor=\color{coderule},
    backgroundcolor=\color{codebg},
    tabsize=2,
    columns=fullflexible,
    keepspaces=true,
    captionpos=b,
    xleftmargin=1.2em,
    aboveskip=0.9em,
    belowskip=0.9em,
}

\newcommand{\code}[1]{\texttt{\small #1}}
\newcommand{\fonteslide}[1]{{\footnotesize\itshape\color{gray} Fonte: slide \emph{L04}, pag.~#1.\par}}
\newcommand{\fontelibro}[1]{{\footnotesize\itshape\color{gray} Fonte: \emph{Thinking in C++} (Eckel), Vol.~2, #1.\par}}
\newcommand{\fonteprove}[1]{{\footnotesize\itshape\color{gray} Misurato in \code{prove/strategy/#1} (compilato ed eseguito).\par}}

\graphicspath{{./img/}}

%==============================================================
\begin{document}

\begin{titlepage}
    \centering
    \vspace*{2.0cm}
    {\Huge\sffamily\bfseries Strategy Pattern\par}
    \vspace{0.35cm}
    {\LARGE\sffamily L'infrastruttura dei linguaggi\par}
    \vspace{0.9cm}
    {\large\sffamily Nested class, inner class, local class\\
    e i loro equivalenti\par}
    \vspace{1.0cm}
    {\LARGE\sffamily Java \quad$\cdot$\quad C++ \quad$\cdot$\quad Python \quad$\cdot$\quad C\par}
    \vspace{1.3cm}
    {\Large Appunti di Ingegneria del Software\par}
    \vspace{1.3cm}
    {\footnotesize\itshape Il pattern in se' e' dato per noto: qui interessa
    \textbf{con quali meccanismi del linguaggio} lo si costruisce, e cosa
    cambia nel comportamento a seconda della forma scelta.\\[6pt]
    Ogni affermazione tecnica di questo documento e' stata
    \textbf{misurata} compilando ed eseguendo
    (g++ 15.2.1, gcc 15.2.1, JDK 21, Python 3.13):\\
    i sorgenti sono in \code{prove/strategy/}.\par}
    \vfill
    {\large Anno 3\par}
\end{titlepage}

\tableofcontents
\newpage

%==============================================================
\chapter{Il problema infrastrutturale}\label{ch:infra}
%==============================================================

\section{Cosa serve, meccanicamente, per fare Strategy}

Il pattern richiede due sole capacita' dal linguaggio:

\begin{enumerate}
    \item \textbf{memorizzare un riferimento a del codice} in una
          struttura dati;
    \item \textbf{sostituirlo} con un altro riferimento.
\end{enumerate}

Tutto il resto (interfacce, classi astratte, ereditarieta') e' un
\emph{modo} di ottenere queste due cose, non un requisito. I meccanismi
disponibili sono quattro: \textbf{puntatori a funzione} nativi,
\textbf{first-class function}, \textbf{classi o istanze di classi},
\textbf{reflection}.

\mprop{Perche' le slide del corso parlano di late binding}{
    Le slide aprono con una motivazione che spiega tutto il resto:
    \begin{quote}\itshape
    ``Here we reflect substantially on how the strategy pattern is
    achieved in Java. This is at the core of polymorphism via overriding
    and late binding. \textbf{Indeed, this is also because we do not use
    pointers to functions.}''
    \end{quote}
    Java non ha puntatori a funzione: deve ottenere lo stesso risultato con
    \textbf{oggetti} e \textbf{late binding}. E' per questo che le slide,
    prima di arrivare al pattern, dedicano trenta pagine a
    \emph{nested class}, \emph{inner class}, \emph{local class},
    \emph{classi anonime} e \emph{interfacce}: sono l'infrastruttura con
    cui, in Java, si impacchetta una strategia.
}

\fonteslide{4}

\section{Il punto di partenza: la frase-chiave delle slide}

\begin{quote}\itshape
``In Java and in C++ it's possible to place a class definition within
another class definition. \textbf{In C++ this is just related to
visibility and naming. In Java the matter is deeper}; we can distinguish
between \emph{nested classes}, \emph{inner classes} and \emph{local
classes}.''
\end{quote}

\fonteslide{6}

Questo documento prende quella frase e la verifica in tutti e quattro i
linguaggi, misurando che cosa cambia davvero.

\section{La mappa: quale forma esiste in quale linguaggio}

\begin{center}
\renewcommand{\arraystretch}{1.35}
\setlength{\tabcolsep}{4pt}
\footnotesize
\begin{tabular}{|>{\raggedright\arraybackslash}p{3.1cm}|>{\centering\arraybackslash}p{2.5cm}|>{\centering\arraybackslash}p{2.5cm}|>{\centering\arraybackslash}p{2.5cm}|>{\centering\arraybackslash}p{2.2cm}|}
\hline
 & \textbf{Java} & \textbf{C++} & \textbf{Python} & \textbf{C} \\
\hline
classe annidata (namespacing) & si' (\code{static}) & si' & si' & \textbf{no}: il nome resta globale \\
\hline
legame implicito con l'istanza esterna & \textbf{si'} (inner class) & no & no & no \\
\hline
classe locale in una funzione & si' & si' & si' & solo la \code{struct}, senza metodi \\
\hline
la classe locale cattura le variabili locali & si' (\emph{effectively final}) & \textbf{NO} & \textbf{si'} (closure) & no \\
\hline
classe anonima / senza nome & si' & lambda & \code{type()} o lambda & no \\
\hline
puntatori a funzione & no & si' & (funzioni first-class) & \textbf{si'} \\
\hline
\end{tabular}
\end{center}

\nt{
    Le due caselle piu' sorprendenti, entrambe verificate col compilatore:
    \begin{itemize}
        \item in \textbf{C++} una classe locale \textbf{non puo'} leggere
              le variabili automatiche della funzione che la contiene ---
              cosa che in Java e Python si fa normalmente;
        \item in \textbf{C} annidare una \code{struct} \textbf{non crea
              alcuno scope}: il nome e' visibile anche fuori (in C++ la
              stessa riga e' un errore).
    \end{itemize}
}

\section{L'esempio comune: \texttt{SimpleCalc}}

Per confrontare i quattro linguaggi a parita' di problema, tutti gli
esempi implementano la stessa \textbf{calcolatrice} delle slide
(\code{SimpleCalc}, pagg.~37--50), resa \emph{headless} per poter girare
senza display grafico:

\begin{itemize}
    \item \textbf{Strategy} $=$ \code{Operazione}, con un metodo
          \code{applica(a, b)};
    \item \textbf{Context} $=$ \code{Calcolatrice}, che tiene un
          accumulatore e l'operazione corrente, e vi \emph{delega};
    \item \textbf{ConcreteStrategy} $=$ le singole operazioni,
          implementate ogni volta con una forma di classe diversa.
\end{itemize}

Il calcolo di prova e' sempre lo stesso: \code{20 \(\langle\)op\(\rangle\)
6}, cosi' i risultati sono confrontabili a colpo d'occhio fra linguaggi.

%==============================================================
\chapter{Java: l'interfaccia come tipo Strategy}\label{ch:interfacce}
%==============================================================

Prima di scegliere \emph{dove} mettere l'implementazione (nested, inner,
local, anonima --- capitolo~\ref{ch:java}), bisogna dichiarare il
\textbf{tipo} della strategia. In Java quel tipo e' quasi sempre
un'\textbf{interfaccia}, ed e' il costrutto a cui le slide dedicano piu'
spazio (pagg.~13--36).

\section{Dichiarazione}

\dfn{Interface (slide 13)}{
    Una dichiarazione di interfaccia introduce un \textbf{nuovo reference
    type} i cui membri sono classi, interfacce, costanti e metodi. Questo
    tipo \textbf{non ha variabili di istanza} e tipicamente dichiara uno o
    piu' \textbf{metodi astratti}; classi altrimenti \emph{non
    imparentate} possono implementarla fornendone le implementazioni.
    \\[3pt]
    Le interfacce \textbf{non possono essere istanziate}: possono solo
    essere implementate da classi o estese da altre interfacce.
    Contengono solo costanti, firme di metodi, default method, static
    method e tipi annidati.
}

\fonteslide{13}

\begin{lstlisting}[language=Java]
public interface Bicycle {
    void changeCadence(int newValue);
    void changeGear(int newValue);
    void speedUp(int increment);
    void applyBrakes(int decrement);
}
\end{lstlisting}

\fonteslide{14}

\mprop{Perche' Java ammette l'ereditarieta' multipla di interfacce}{
    Le slide danno la ragione precisa: le interfacce
    \emph{``support single and multiple inheritance, since \textbf{the
    absence of instance data eliminates the ambiguities} of, say, C++''}.
    \\[3pt]
    Non avendo stato di istanza, non esiste il problema dei
    \textbf{dati duplicati} che in C++ richiede l'ereditarieta' virtuale.
    Resta solo il problema del \emph{comportamento} duplicato, risolto
    dalle due regole della sezione~\ref{sec:regole}.
}

\fonteslide{14}

\begin{Example}{Verificato: un'interfaccia ha davvero solo costanti}{}
Interrogando a runtime i campi dichiarati in un'interfaccia con costanti:
\begin{lstlisting}
campi dichiarati in un'interfaccia:
  E  ->  static=true final=true public=true
\end{lstlisting}
Tutte le costanti sono \textbf{implicitamente} \code{public static
final}; tutti i metodi abstract, default e static sono
\textbf{implicitamente} \code{public}. Nessuno stato di istanza, come
promesso.
\end{Example}

\fonteslide{15}
\fonteprove{NL5\_Java\_Interfacce.java}

\section{Implementazione}

\begin{lstlisting}[language=Java]
class BianchiBicycle implements Bicycle {
  int cadence = 0, speed = 0, gear = 1;
  // Il compilatore RICHIEDE che changeCadence, changeGear, speedUp e
  // applyBrakes siano tutti implementati: la compilazione FALLISCE se
  // uno di essi manca.
  void changeCadence(int newValue) { cadence = newValue; }
  // ...
}
\end{lstlisting}

\fonteslide{16}

\section{Default e static method}

\dfn{Default e static method (slide 15, 29)}{
    Un \textbf{default method} (parola chiave \code{default}) ha un corpo
    e \textbf{puo' essere sovrascritto}: e' il meccanismo che permette
    l'\emph{evoluzione} di un'interfaccia, supportando l'implementazione
    non obbligatoria di tutti i metodi nelle classi.
    \\[3pt]
    Un \textbf{static method} serve per comportamenti generali
    dell'interfaccia: \textbf{non} puo' accedere a dati non-static e
    \textbf{non} puo' essere sovrascritto.
}

\fonteslide{29}

\begin{Example}{Il problema che i default risolvono (esempio Oracle)}{}
Data un'interfaccia gia' pubblicata e usata da altri:
\begin{lstlisting}[language=Java]
public interface DoIt {
   void doSomething(int i, double x);
   int  doSomethingElse(String s);
}
\end{lstlisting}
volendo aggiungere un metodo, \textbf{tutto il codice delle vecchie
implementazioni si romperebbe}, perche' non ha quel metodo. Con
\code{default} il problema sparisce:
\begin{lstlisting}[language=Java]
public interface DoIt {
   void doSomething(int i, double x);
   int  doSomethingElse(String s);
   default boolean didItWork(int i, double x, String s) {
       // Method body
   }
}
\end{lstlisting}
\textbf{Verificato:} una classe scritta \emph{prima} dell'aggiunta di
\code{didItWork} continua a compilare e a funzionare, ereditando il
default.
\end{Example}

\fonteslide{30-32}
\fonteprove{NL5\_Java\_Interfacce.java}

\section{Le due regole di risoluzione dei conflitti}\label{sec:regole}

Default e abstract method si ereditano come metodi di istanza. Ma quando
i supertipi forniscono \textbf{piu' default method con la stessa firma},
il compilatore applica due principi:

\mprop{Le due regole (slide 17)}{
    \begin{enumerate}
        \item i \textbf{metodi di istanza sono preferiti} ai default
              method delle interfacce;
        \item i metodi \textbf{gia' sovrascritti da altri candidati sono
              ignorati} --- circostanza che si presenta quando i supertipi
              condividono un antenato comune.
    \end{enumerate}
}

\fonteslide{17}

\subsection{Regola 1: l'istanza vince}

\begin{lstlisting}[language=Java]
public class Horse {
  public String identifyMyself() { return "I am a horse."; }
}
public interface Flyer {
  default public String identifyMyself() { return "I am able to fly."; }
}
public interface Mythical {
  default public String identifyMyself() {
      return "I am a mythical creature."; }
}
public class Pegasus extends Horse implements Flyer, Mythical { }
\end{lstlisting}

\textbf{Output verificato:} \code{Pegasus.identifyMyself()} restituisce
\emph{``I am a horse.''} --- il metodo di istanza ereditato dalla
\emph{classe} \code{Horse} vince sui default delle due interfacce.

\fonteslide{18-19}

\subsection{Regola 2: i gia' sovrascritti sono ignorati}

\begin{lstlisting}[language=Java]
public interface Animal {
    default public String identifyMyself() { return "I am an animal."; }
}
public interface EggLayer extends Animal {
    default public String identifyMyself() { return "I am able to lay eggs."; }
}
public interface FireBreather extends Animal { }      // NON fa override

public class Dragon implements EggLayer, FireBreather { }
\end{lstlisting}

\textbf{Output verificato:} \emph{``I am able to lay eggs.''} ---
\code{EggLayer} ha gia' sovrascritto il default di \code{Animal}, quindi
quello di \code{Animal} (che \code{FireBreather} porterebbe) viene
\textbf{ignorato}. Nessuna ambiguita', nessun override necessario.

\fonteslide{20-21}

\subsection{Il conflitto vero: due default indipendenti}

Quando due default \textbf{indipendenti} confliggono (o un default
confligge con un abstract), il compilatore \textbf{produce un errore} e
si e' obbligati a fare l'override esplicito.

\begin{lstlisting}[language=Java]
public interface OperateCar {
    default public int startEngine(EncryptedKey key) { /* impl */ }
}
public interface FlyCar {
    default public int startEngine(EncryptedKey key) { /* impl */ }
}
public class FlyingCar implements OperateCar, FlyCar { }   // ERRORE
\end{lstlisting}

\textbf{Errore reale prodotto dal compilatore:}

\begin{lstlisting}
error: types OperateCar and FlyCar are incompatible;
  class FlyingCar inherits unrelated defaults for
  startEngine(EncryptedKey) from types OperateCar and FlyCar
\end{lstlisting}

La correzione, con la sintassi \code{Interfaccia.super.metodo()}:

\begin{lstlisting}[language=Java]
public class FlyingCar implements OperateCar, FlyCar {
    public int startEngine(EncryptedKey key) {
        FlyCar.super.startEngine(key);
        OperateCar.super.startEngine(key);
    }
}
\end{lstlisting}

\nt{
    Il nome che precede \code{super} deve riferirsi a una
    \textbf{superinterfaccia diretta} che definisce o eredita un default
    per quel metodo. Questa forma di invocazione non e' limitata al caso
    del conflitto: \code{super} si puo' usare per invocare un default
    method sia in classi sia in interfacce.
}

\fonteslide{22-25}
\fonteprove{NL5\_Java\_Interfacce.java}

\subsection{Un metodo di istanza ereditato soddisfa un abstract}

\begin{lstlisting}[language=Java]
public interface Mammal { String identifyMyself(); }   // ASTRATTO
public class Horse { public String identifyMyself() { return "I am a horse."; } }

public class Mustang extends Horse implements Mammal { }   // non implementa nulla!
\end{lstlisting}

\textbf{Output verificato:} \emph{``I am a horse.''} --- \code{Mustang}
eredita \code{identifyMyself()} da \code{Horse}, e quel metodo concreto
\textbf{soddisfa} il contratto astratto di \code{Mammal}.

\fonteslide{26-27}

\subsection{Il riepilogo delle slide}

\begin{center}
\renewcommand{\arraystretch}{1.4}
\footnotesize
\begin{tabular}{|>{\raggedright\arraybackslash}p{4.0cm}|>{\centering\arraybackslash}p{4.0cm}|>{\centering\arraybackslash}p{4.0cm}|}
\hline
 & \textbf{Metodo di istanza della superclasse} & \textbf{Metodo static della superclasse} \\
\hline
\textbf{Metodo di istanza della sottoclasse} & \emph{Overrides} & errore di compilazione \\
\hline
\textbf{Metodo static della sottoclasse} & errore di compilazione & \emph{Hides} \\
\hline
\end{tabular}
\end{center}

\fonteslide{28}

\section{Il listener: l'interfaccia a un solo metodo}

Le slide arrivano allo Strategy attraverso gli \textbf{event listener} di
Swing. Un listener \emph{e'} un'interfaccia con un solo metodo:

\begin{lstlisting}[language=Java]
public interface ActionListener extends EventListener {
    public void actionPerformed(ActionEvent e);
}
\end{lstlisting}

\fonteslide{36}

\mprop{ActionListener \emph{e'} una Strategy}{
    \begin{itemize}
        \item l'\textbf{interfaccia} \code{ActionListener} e' il tipo
              Strategy: incapsula ``cosa fare quando succede l'evento'';
        \item il \code{JButton} e' il \textbf{Context}: non sa cosa fara'
              il listener, lo invoca e basta;
        \item \code{addActionListener()} e' l'\textbf{iniezione} della
              strategia.
    \end{itemize}
    Le slide notano anche il vantaggio pratico: piu' componenti grafici
    possono condividere \textbf{lo stesso} listener, se fanno la stessa
    cosa --- che e' esattamente il riuso promesso dallo Strategy.
}

\fonteslide{34-35}

\section{Functional interface: il ponte verso le lambda}

\begin{Prop}{Verificato: una lambda richiede UN SOLO metodo astratto}{}
    Contando a runtime i metodi astratti dichiarati:
\begin{lstlisting}
MioActionListener    metodi astratti: 1  -> e' una FUNCTIONAL INTERFACE
Bicycle              metodi astratti: 4  -> NON e' functional
\end{lstlisting}
    Provando a implementare \code{Bicycle} con una lambda:
\begin{lstlisting}
error: incompatible types: Bicycle is not a functional interface
    multiple non-overriding abstract methods found in interface Bicycle
\end{lstlisting}
    E' per questo che i \textbf{default method contano} anche qui:
    aggiungere default method a un'interfaccia \textbf{non} le fa perdere
    lo status di functional interface, perche' i default non sono
    astratti. Un'interfaccia puo' avere venti default method e restare
    implementabile con una lambda.
\end{Prop}

\fonteprove{NL5\_Java\_Interfacce.java}

%==============================================================
\chapter{Java: quattro forme piu' la lambda}\label{ch:java}
%==============================================================

\section{Le tre categorie delle slide}

\dfn{Nested class}{
    Una classe dentro un'altra, dichiarata \code{static}. Serve per
    \textbf{naming e packaging}: permette di \emph{nascondere l'esistenza}
    di una classe al mondo esterno e di raggruppare classi che
    appartengono logicamente insieme. Il nome e' \textbf{locale alla
    classe che la racchiude}.
    \\[3pt]
    \textbf{Vincolo:} le dichiarazioni in una nested class possono usare
    \textbf{solo i membri \code{static}} della classe esterna.
}

\fonteslide{7}

\dfn{Inner class}{
    Una classe dentro un'altra \textbf{senza} \code{static}. In Java
    ``dichiarare semplicemente una classe dentro il corpo di un'altra,
    senza ulteriori qualificazioni, produce una inner class''.
    \\[3pt]
    Ogni istanza ha una \textbf{enclosing instance} e puo' usarne i
    membri --- \emph{anche quelli di istanza}. Non puo' avere membri
    \code{static}, ne' lo stesso nome della classe che la racchiude.
}

\fonteslide{8, 11}

\dfn{Local class}{
    Una classe dichiarata \textbf{dentro il corpo di una funzione}. Utile
    quando servono codice e strutture dati per compiti specifici dentro
    una funzione, senza renderli disponibili all'esterno: e' un mezzo per
    imporre \textbf{incapsulamento a livello di scope di funzione}. Il
    nome e' locale allo scope che la racchiude.
    \\[3pt]
    In Java le local class sono \textbf{a tutti gli effetti inner class},
    non semplici nested: accedono ai membri della classe esterna, alle
    variabili locali \code{final} e ai parametri.
}

\fonteslide{12}

\begin{Example}{L'esempio canonico delle slide}{}
\begin{lstlisting}[language=Java]
public class X {
  int instanceVar;           // appartiene a un'istanza di X
  static int staticVar;      // appartiene alla classe X

  private class Inner {                      // INNER (non static)
    int y;
    Inner() { y = instanceVar + staticVar; } // puo' usare ENTRAMBI
  }

  private static class Nested {              // NESTED (static)
    int x;
    Nested() {
      x = staticVar;
      // x += instanceVar;   // Errore! nessuna istanza di X disponibile
    }
  }
}
\end{lstlisting}
La differenza sta in \emph{una sola parola chiave}, ma cambia la natura
della classe: \code{Nested} e' una classe come le altre, \code{Inner} non
esiste senza un oggetto \code{X} che la ospiti.
\end{Example}

\fonteslide{9}

\section{La sintassi di creazione}

\begin{Prop}{\texttt{esterno.new Inner()}}{}
    Una nested class si istanzia come una classe normale, col nome
    qualificato. Una \textbf{inner} class richiede una sintassi speciale,
    perche' serve \emph{un'istanza dell'esterno}:
\begin{lstlisting}[language=Java]
// nested (static): il nome e' qualificato dalla CLASSE
Calcolatrice.Somma s = new Calcolatrice.Somma();

// inner: il nome e' qualificato da un'ISTANZA
NL1_Java_SimpleCalc esterno = new NL1_Java_SimpleCalc();
Operazione inner = esterno.new MoltiplicaPerCampo();
\end{lstlisting}
    Provare a scrivere \code{new MoltiplicaPerCampo()} da un contesto
    \code{static} e' un errore di compilazione: non c'e' nessuna
    enclosing instance da agganciare.
\end{Prop}

\section{Le tre classi dello Strategy in Java}

\subsection{Strategy e Context}

\begin{lstlisting}[language=Java]
// STRATEGY: un'interfaccia ANNIDATA (le interfacce annidate sono
// sempre implicitamente static)
interface Operazione {
    int applica(int a, int b);
    default String simbolo() { return "?"; }
}

// CONTEXT: tiene lo stato (l'accumulatore) e delega alla strategia
static class Calcolatrice {
    private int accumulatore = 0;
    private Operazione operazione;             // la STRATEGIA

    void setOperazione(Operazione op) { this.operazione = op; }

    void premiNumero(int n) {                  // i "bottoni" delle slide
        if (operazione == null) { accumulatore = n; return; }
        accumulatore = operazione.applica(accumulatore, n);  // <- delega
    }
    int display() { return accumulatore; }
}
\end{lstlisting}

\subsection{Forma A --- static nested}

\begin{lstlisting}[language=Java]
static int contatoreStatico = 0;        // membro STATIC dell'esterna
private int campoDiIstanza = 100;       // membro di ISTANZA

static class Somma implements Operazione {   // <- static: e' NESTED
    @Override public int applica(int a, int b) {
        contatoreStatico++;             // OK: membro static
        // campoDiIstanza;              // ERRORE: non c'e' un'istanza!
        return a + b;
    }
    @Override public String simbolo() { return "+"; }
}
\end{lstlisting}

\textbf{Perche' questa forma:} la strategia non ha bisogno di conoscere lo
stato di un particolare oggetto \code{Calcolatrice}; \code{static} lo
dichiara esplicitamente ed evita di tenere in vita l'oggetto esterno.
E' la forma da preferire per default.

\subsection{Forma B --- inner class}

\begin{lstlisting}[language=Java]
class MoltiplicaPerCampo implements Operazione {   // NON static: INNER
    @Override public int applica(int a, int b) {
        // legge il campo di ISTANZA dell'oggetto esterno:
        return a * b % campoDiIstanza;
    }
    @Override public String simbolo() { return "*%" + campoDiIstanza; }
}
\end{lstlisting}

\textbf{Perche' questa forma:} e' la scelta delle slide per i bottoni
\code{+} e \code{-} di \code{SimpleCalc}. Il commento della slide 43 e'
esplicito: \emph{``this also has the advantage of allowing the inner class
access to the components defined within the outer class''}. La strategia
deve leggere e scrivere \code{currentCalc} e \code{calcOperation} della
calcolatrice: una inner class lo fa senza che si debba passare nulla.

\begin{Example}{La prova che il legame e' vivo (output reale)}{}
Lo \textbf{stesso} oggetto-strategia, dopo aver modificato un campo
dell'\emph{oggetto esterno}:
\begin{lstlisting}
campoDiIstanza=100 -> 20*6 % 100 = 20
campoDiIstanza=7   -> 20*6 % 7   = 1
\end{lstlisting}
E creando l'inner class da \emph{un'altra} istanza esterna:
\begin{lstlisting}
inner di un'ALTRA istanza (campo=1000): 120
\end{lstlisting}
Non si tratta di una copia del valore fatta alla costruzione: la
strategia \textbf{legge lo stato dell'enclosing instance al momento della
chiamata}. Due istanze esterne danno due ``mondi'' separati.
\end{Example}

\fonteprove{NL1\_Java\_SimpleCalc.java}

\subsection{Forma C --- local class}

\begin{lstlisting}[language=Java]
static Operazione creaSottrazioneLimitata(int minimo) {
    // 'minimo' e' effectively final: catturabile
    class SottrazioneLimitata implements Operazione {   // visibile SOLO qui
        @Override public int applica(int a, int b) {
            return Math.max(a - b, minimo);   // <- CATTURA 'minimo'
        }
        @Override public String simbolo() { return "-(>=" + minimo + ")"; }
    }
    return new SottrazioneLimitata();
    // fuori dal metodo il NOME 'SottrazioneLimitata' non esiste piu'
}
\end{lstlisting}

\textbf{Perche' questa forma:} la classe serve solo a costruire
\emph{quel} particolare oggetto, e non ha senso come tipo pubblico. Il
nome resta confinato al metodo (incapsulamento a livello di funzione), ma
il \emph{tipo} continua a esistere: l'oggetto restituito e' perfettamente
utilizzabile fuori, attraverso l'interfaccia.

\mprop{La regola della cattura: \emph{effectively final}}{
    Una local class (e una classe anonima) puo' catturare una variabile
    locale solo se questa non viene mai riassegnata dopo
    l'inizializzazione. Il motivo e' che la cattura avviene
    \textbf{per valore}: il compilatore copia il valore nell'oggetto, e
    permettere riassegnazioni successive creerebbe due copie divergenti.
    \\[3pt]
    Nota la differenza rispetto ai \emph{campi}: la forma B legge
    \code{campoDiIstanza} \emph{ogni volta}, e infatti lo vede cambiare;
    la forma C ha una copia di \code{minimo} congelata alla creazione.
}

\subsection{Forma D --- classe anonima}

\begin{lstlisting}[language=Java]
static Operazione divisioneSicura() {
    return new Operazione() {              // <- classe SENZA nome
        @Override public int applica(int a, int b) {
            return b == 0 ? a : a / b;
        }
        @Override public String simbolo() { return "/"; }
    };
}
\end{lstlisting}

\textbf{Perche' questa forma:} le slide la introducono con una motivazione
di ingegneria del software, non di sintassi. Quando \emph{ogni} bottone ha
un comportamento diverso, avere una classe con nome per ciascuno
provocherebbe una \emph{``proliferation of standard classes, making the
code not readable''}. Con le classi anonime \emph{``the behaviour
associated with a button is placed right where such button is
instantiated''}.

\fonteslide{47}

\wc{La classe anonima e' ``un pezzo di codice passato come parametro''}{
    Le slide insistono su questo punto: \emph{``we get the impression that
    a piece of code could be the parameter of the call of the method
    \code{addActionListener}. \textbf{But it is not so}: it is a fully
    fledged (anonymous) class that is created, and an object of it is then
    instantiated and passed as a parameter''}.
    \\[3pt]
    E' la conferma della slide 4: Java non passa funzioni, passa
    \textbf{oggetti}. La sintassi concisa maschera un oggetto vero.
}

\fonteslide{49}

\subsection{Forma E --- lambda}

\begin{lstlisting}[language=Java]
c.setOperazione((a, b) -> a > b ? a : b);   // lambda
c.setOperazione(Math::sqrt);                 // method reference
\end{lstlisting}

Possibile perche' \code{Operazione} ha un solo metodo astratto
(\emph{functional interface} de facto).

\section{La misura decisiva: i nomi generati dal compilatore}

Le cinque forme non sono varianti stilistiche: producono \textbf{artefatti
diversi}. Chiedendo a runtime il nome della classe di ciascuna strategia,
si ottiene (output reale):

\begin{lstlisting}
A static nested -> NL1_Java_SimpleCalc$Somma
B inner         -> NL1_Java_SimpleCalc$MoltiplicaPerCampo
C local         -> NL1_Java_SimpleCalc$1SottrazioneLimitata
D anonima       -> NL1_Java_SimpleCalc$1
E lambda        -> NL1_Java_SimpleCalc$$Lambda/0x000000003f041000
\end{lstlisting}

E i file effettivamente prodotti da \code{javac}:

\begin{lstlisting}
NL1_Java_SimpleCalc$1.class                      <- anonima
NL1_Java_SimpleCalc$1SottrazioneLimitata.class   <- local
NL1_Java_SimpleCalc$Calcolatrice.class
NL1_Java_SimpleCalc$MoltiplicaPerCampo.class     <- inner
NL1_Java_SimpleCalc$Operazione.class
NL1_Java_SimpleCalc$Somma.class                  <- nested
NL1_Java_SimpleCalc.class
\end{lstlisting}

\mprop{Cosa insegna questo elenco}{
    \begin{itemize}
        \item nested e inner producono un \code{.class} col
              \textbf{proprio nome} (\code{\$Somma},
              \code{\$MoltiplicaPerCampo}): sono tipi a pieno titolo,
              distinguibili;
        \item la \textbf{local class} ha un \textbf{numero davanti al
              nome} (\code{\$1SottrazioneLimitata}): il numero serve
              perche' due metodi diversi potrebbero dichiarare due local
              class omonime;
        \item la \textbf{anonima} ha \textbf{solo un numero}
              (\code{\$1}): non ha un nome da conservare;
        \item la \textbf{lambda} \textbf{non ha alcun file
              \code{.class}}. Il suo nome contiene un indirizzo
              (\code{/0x\dots}) perche' la classe e' \textbf{generata a
              runtime} dalla JVM (\code{invokedynamic} $+$
              \code{LambdaMetafactory}).
    \end{itemize}
    Quindi una lambda \textbf{non e'} zucchero sintattico per una classe
    anonima: sono due meccanismi diversi con la stessa apparenza.
}

\fonteprove{NL1\_Java\_SimpleCalc.java}

%==============================================================
\chapter{C++: solo nested, piu' le lambda}\label{ch:cpp}
%==============================================================

\section{L'equivalente dell'interfaccia: la classe astratta pura}

In C++ non esiste \code{interface}: si usa una \textbf{classe astratta
pura} --- soli metodi virtuali puri, nessun dato di istanza --- piu' un
distruttore virtuale.

\begin{lstlisting}[language=C++]
struct Bicycle {                      // "interfaccia" in C++
    virtual void changeCadence(int) = 0;
    virtual void changeGear(int) = 0;
    virtual void speedUp(int) = 0;
    virtual void applyBrakes(int) = 0;
    virtual ~Bicycle() = default;     // obbligatorio: si distrugge via base*
};
\end{lstlisting}

E l'equivalente di un \textbf{default method} e' un metodo virtuale
\textbf{con corpo}. Ma qui la corrispondenza si rompe, e in modo
istruttivo.

\wc{Un metodo virtuale con corpo e' l'equivalente di un default method}{
    \textbf{Non del tutto.} Un default method Java non aggiunge ne' stato
    ne' un sotto-oggetto alla classe che implementa l'interfaccia. Un
    metodo virtuale C++ porta con se' una \textbf{base vera}, con vtable
    e sotto-oggetto --- e quindi riapre il \emph{diamond problem} che le
    interfacce Java evitano per costruzione.
}

\subsection{I tre casi delle slide, tradotti}

\begin{center}
\renewcommand{\arraystretch}{1.35}
\setlength{\tabcolsep}{4pt}
\footnotesize
\begin{tabular}{|>{\raggedright\arraybackslash}p{2.3cm}|>{\raggedright\arraybackslash}p{5.0cm}|>{\raggedright\arraybackslash}p{5.0cm}|}
\hline
\textbf{Caso} & \textbf{Java} & \textbf{C++} \\
\hline
\emph{Pegasus} (istanza vs default) & compila da solo: vince \code{Horse} (regola 1) & \textbf{ambiguo}: \emph{``request for member 'identifyMyself' is ambiguous''}. Va risolto a mano \\
\hline
\emph{Dragon} (antenato comune) & compila da solo: vince \code{EggLayer} (regola 2) & serve l'\textbf{ereditarieta' virtuale}; poi funziona per \textbf{dominanza} \\
\hline
\emph{FlyingCar} (default indipendenti) & errore, override obbligatorio, \code{FlyCar.super.f()} & \textbf{identico}, ma la sintassi e' \code{FlyCar::f()} \\
\hline
\end{tabular}
\end{center}

\begin{Prop}{Il caso Dragon e' l'analogo esatto della regola 2}{}
    In C++ il caso \code{Dragon} si scrive cosi':
\begin{lstlisting}[language=C++]
struct Animal       { virtual std::string identifyMyself() const
                      { return "I am an animal."; } };
struct EggLayer     : virtual Animal {          // <- virtual!
    std::string identifyMyself() const override
    { return "I am able to lay eggs."; } };
struct FireBreather : virtual Animal { };
struct Dragon : EggLayer, FireBreather { };     // nessun override
\end{lstlisting}
    Funziona (output verificato: \emph{``I am able to lay eggs.''})
    grazie alla \textbf{dominanza}: con \code{virtual} c'e' un solo
    sotto-oggetto \code{Animal}, ed \code{EggLayer} lo domina.
    \\[3pt]
    E' \emph{la stessa idea} della regola 2 di Java (``i metodi gia'
    sovrascritti sono ignorati''), ottenuta pero' con un meccanismo del
    tutto diverso --- e che in Java non serve, perche' le interfacce non
    portano sotto-oggetti da condividere.
\end{Prop}

\begin{Example}{La misura che spiega tutto: \code{sizeof}}{}
\begin{lstlisting}
sizeof(Pegasus) = 24 byte     (tre basi polimorfiche -> TRE vptr)
\end{lstlisting}
In Java implementare dieci interfacce \textbf{non aggiunge nulla}
all'oggetto: le interfacce non hanno stato di istanza. In C++ ogni base
polimorfica costa un vptr. E' la conferma quantitativa della frase della
slide 14 (\emph{``the absence of instance data eliminates the ambiguities
of, say, C++''}): non e' solo una questione di ambiguita' logica, e'
anche una questione di \textbf{layout in memoria}.
\end{Example}

\fonteprove{NL6\_interfacce\_altri\_linguaggi.cpp}

\section{``In C++ this is just related to visibility and naming''}

La frase della slide 6 e' esatta, e ha tre conseguenze precise, tutte
verificate.

\mprop{Conseguenza 1: non esiste l'inner class}{
    In C++ una classe annidata e' \textbf{sempre} l'equivalente di una
    \code{static} nested class di Java: \textbf{non} ha alcun legame
    implicito con un'istanza della classe esterna.
}

\begin{lstlisting}[language=C++]
struct E {
    int v = 1;
    struct N {
        int f() { return v; }   // ERRORE: invalid use of non-static
    };                          //         data member 'E::v'
};
\end{lstlisting}

Per ottenere il comportamento di una inner class Java bisogna
\textbf{passare esplicitamente} un puntatore all'oggetto esterno --- che
e' esattamente cio' che il compilatore Java genera di nascosto.

\mprop{Conseguenza 2: ma l'accesso ai privati c'e'}{
    Dal C++11 una classe annidata ha accesso ai membri \textbf{privati}
    della classe che la racchiude (ha gli stessi diritti di un membro).
    Verificato: la nested \code{Somma} incrementa il
    \code{contatoreStatico\_} privato dell'enclosing.
}

\mprop{Conseguenza 3: lo scope, invece, c'e' eccome}{
    Il nome di una classe annidata e' \textbf{qualificato}:
    \code{Calcolatrice::Somma}. Non e' visibile come \code{Somma} nudo
    all'esterno. (Si vedra' nel capitolo sul C che questo, per le
    \code{struct}, e' una differenza C/C++ sostanziale.)
}

\section{Le tre classi dello Strategy in C++}

\subsection{Strategy e Context, con le strategie annidate}

\begin{lstlisting}[language=C++]
// STRATEGY: in C++ non esiste 'interface' -> classe astratta pura
struct Operazione {
    virtual int applica(int a, int b) const = 0;
    virtual std::string simbolo() const { return "?"; }
    virtual ~Operazione() = default;      // sempre, se si distrugge via base*
};

class Calcolatrice {                       // CONTEXT
    int accumulatore_ = 0;
    const Operazione* operazione_ = nullptr;
    static int contatoreStatico_;          // membro STATIC privato
    int campoDiIstanza_ = 100;             // membro di ISTANZA privato

public:
    // --- FORMA A: nested class ---
    struct Somma : Operazione {
        int applica(int a, int b) const override {
            contatoreStatico_++;           // OK: privato static dell'enclosing
            // return a + b + campoDiIstanza_;  // ERRORE: nessuna istanza!
            return a + b;
        }
        std::string simbolo() const override { return "+"; }
    };
    // ...
};
int Calcolatrice::contatoreStatico_ = 0;   // definizione fuori dalla classe
\end{lstlisting}

\subsection{Forma B --- l'equivalente dell'inner class, a mano}

\begin{lstlisting}[language=C++]
class MoltiplicaPerCampo : public Operazione {
    const Calcolatrice* esterno_;          // <- l'"enclosing instance", A MANO
public:
    explicit MoltiplicaPerCampo(const Calcolatrice* e) : esterno_(e) {}
    int applica(int a, int b) const override {
        return a * b % esterno_->campoDiIstanza_;   // accede al privato
    }
};

// factory che imita la sintassi Java 'esterno.new Inner()'
std::unique_ptr<Operazione> creaMoltiplicaPerCampo() const {
    return std::make_unique<MoltiplicaPerCampo>(this);
}
\end{lstlisting}

\textbf{Perche' questa struttura:} rende \emph{visibile nel codice} cio'
che Java nasconde. Il costo dell'``inner class'' Java c'e' anche in Java
--- un campo puntatore in piu' per oggetto --- solo che non si vede. In
C++ lo si paga esplicitamente, e infatti si misura:

\begin{lstlisting}
sizeof(Calcolatrice::Somma)        = 8    (solo il vptr)
sizeof(MoltiplicaPerCampo)         = 16   (vptr + puntatore all'esterno)
\end{lstlisting}

\fonteprove{NL2\_cpp\_SimpleCalc.cpp}

\subsection{Forma C --- local class, e i suoi due vincoli}

Le slide dicono che ``sia Java sia C++ permettono di definire una classe
dentro il corpo di una funzione''. E' vero, ma in C++ con due limiti che
Java non ha, entrambi verificati con messaggi d'errore reali.

\begin{wconc}{Una local class C++ puo' usare le variabili locali della funzione}{}
    \textbf{No.} Puo' usare solo variabili \code{static}, globali, enum e
    tipi. Le variabili \textbf{automatiche} sono vietate:
\begin{lstlisting}[language=C++]
void g() {
    int automatica = 1;
    class L {
    public:
        int f() { return automatica; }   // ERRORE
    };
}
\end{lstlisting}
    Messaggio reale del compilatore:
\begin{lstlisting}
error: use of local variable with automatic storage
       from containing function
\end{lstlisting}
    \textbf{Il motivo:} una local class e' un \emph{tipo}, e i suoi metodi
    sono funzioni ordinarie; non esiste alcun meccanismo che copi le
    variabili locali dentro l'oggetto. In Java quel meccanismo esiste (la
    cattura per valore delle \emph{effectively final}); in C++ e' stato
    aggiunto solo nel 2011, e si chiama \textbf{lambda}.
\end{wconc}

\begin{wconc}{Una local class C++ puo' avere membri dati \texttt{static}}{}
    \textbf{No.} Messaggio reale:
\begin{lstlisting}
error: local class 'class g()::L' shall not have
       static data member 'int g()::L::m'
\end{lstlisting}
    Un membro static richiederebbe una definizione a livello di
    namespace, che per una classe locale non e' esprimibile.
\end{wconc}

Nell'esempio eseguibile, per aggirare il primo vincolo, il valore
``catturato'' viene passato tramite una variabile \code{static} locale ---
e il commento nel codice segnala che si tratta di un \emph{workaround},
non della soluzione idiomatica:

\begin{lstlisting}[language=C++]
std::unique_ptr<Operazione> creaSottrazioneLimitata(int minimo) {
    static int minimoStatico = 0;      // workaround: non 'minimo' diretto
    minimoStatico = minimo;

    class SottrazioneLimitata : public Operazione {
    public:
        int applica(int a, int b) const override {
            // return std::max(a - b, minimo);   // *** ERRORE ***
            int m = minimoStatico;               // ...quindi la statica
            return (a - b) > m ? (a - b) : m;
        }
    };
    return std::make_unique<SottrazioneLimitata>();
}
\end{lstlisting}

\nt{
    Questo workaround ha un difetto grave che vale la pena vedere: la
    variabile \code{static} e' \textbf{condivisa fra tutte} le chiamate.
    Creando due strategie con minimi diversi, la seconda
    \emph{sovrascriverebbe} il valore visto dalla prima. E' precisamente
    il problema che le closure risolvono, e il motivo per cui in C++
    moderno per questo scopo si usa una lambda, non una local class.
}

\subsection{Forma D --- la lambda: la vera classe locale del C++}

\begin{lstlisting}[language=C++]
int minimo = 5;                              // variabile AUTOMATICA
cf.setOperazione([minimo](int a, int b) {    // <- catturata!
    int r = a - b;
    return r > minimo ? r : minimo;
}, "-(>=5) lambda");
\end{lstlisting}

\mprop{Cos'e' davvero una lambda}{
    Il compilatore genera un \textbf{tipo senza nome} (\emph{closure
    type}), \textbf{locale} allo scope in cui la lambda appare, con un
    \code{operator()} e un membro per ogni variabile catturata. E' quindi
    letteralmente una \textbf{classe locale anonima con cattura}: risolve
    in un colpo il limite della local class e fornisce l'equivalente della
    classe anonima di Java.
}

\fontelibro{cap.~10, ``Strategy''}

Le dimensioni misurate rendono concreto il confronto:

\begin{center}
\renewcommand{\arraystretch}{1.3}
\footnotesize
\begin{tabular}{|>{\raggedright\arraybackslash}p{6.4cm}|>{\centering\arraybackslash}p{1.8cm}|>{\raggedright\arraybackslash}p{4.0cm}|}
\hline
\textbf{Forma} & \textbf{sizeof} & \textbf{Composizione} \\
\hline
nested class con un metodo virtuale & 8 & solo il vptr \\
\hline
``inner class'' con puntatore all'esterno & 16 & vptr $+$ puntatore \\
\hline
lambda che cattura un \code{int} & \textbf{4} & solo la cattura, \textbf{nessun vptr} \\
\hline
lambda senza cattura & \textbf{1} & oggetto vuoto \\
\hline
\code{std::function<int(int,int)>} & 32 & type erasure \\
\hline
\end{tabular}
\end{center}

\nt{
    Una lambda e' \emph{piu' piccola} di una strategia con funzione
    virtuale, perche' non ha vptr: il tipo e' noto staticamente. Il prezzo
    e' che ogni lambda ha un \textbf{tipo diverso}, quindi per metterle in
    una variabile comune serve \code{std::function} --- che reintroduce
    l'indirezione (e i 32 byte). E' lo stesso trade-off
    monomorfizzazione/erasure incontrato nei template.
}

\fonteprove{NL2\_cpp\_SimpleCalc.cpp}

%==============================================================
\chapter{Python: un solo costrutto, tre comportamenti}\label{ch:python}
%==============================================================

\section{L'equivalente dell'interfaccia: tre possibilita'}

Python offre \textbf{tre} modi di dichiarare il tipo di una strategia,
che coprono uno spettro che in Java non esiste.

\begin{center}
\renewcommand{\arraystretch}{1.35}
\footnotesize
\begin{tabular}{|>{\raggedright\arraybackslash}p{3.4cm}|>{\raggedright\arraybackslash}p{4.4cm}|>{\raggedright\arraybackslash}p{4.4cm}|}
\hline
\textbf{Meccanismo} & \textbf{Tipo di controllo} & \textbf{Quando fallisce} \\
\hline
\code{abc.ABC} $+$ \code{@abstractmethod} & \textbf{nominale}: serve ereditare esplicitamente & a \textbf{runtime}, all'istanziazione \\
\hline
\code{typing.Protocol} & \textbf{strutturale}: basta avere i metodi & a \textbf{compile time}, ma solo se lanci \code{mypy} \\
\hline
niente (duck typing) & nessuno & a runtime, alla prima chiamata mancante \\
\hline
\end{tabular}
\end{center}

\begin{lstlisting}[language=Python]
# 1. NOMINALE: l'equivalente piu' vicino a 'interface' di Java
class Bicycle(ABC):
    @abstractmethod
    def change_cadence(self, v: int) -> None: ...
    def descrivi(self) -> str:              # <- l'equivalente di un DEFAULT
        return f"{type(self).__name__} generica"

# 2. STRUTTURALE: nessuna ereditarieta' richiesta
@runtime_checkable
class Operazione(Protocol):
    def applica(self, a: int, b: int) -> int: ...

class Somma:                                # NON eredita da Operazione!
    def applica(self, a, b): return a + b
\end{lstlisting}

\textbf{Verificato:} \code{issubclass(Somma, Operazione)} restituisce
\code{True} pur non essendoci alcuna parentela dichiarata --- conta la
\emph{struttura}. Ed e' l'opposto esatto di Java, dove serve
\code{implements}.

\fonteprove{NL7\_interfacce\_python.py}

\section{I tre casi delle slide, risolti dall'MRO}

Python non ha ne' le due regole di Java ne' la dominanza del C++: ha
l'\textbf{MRO} (Method Resolution Order, linearizzazione C3). I risultati
spesso coincidono, ma il meccanismo e' diverso --- e le differenze si
vedono.

\begin{center}
\renewcommand{\arraystretch}{1.3}
\footnotesize
\begin{tabular}{|>{\raggedright\arraybackslash}p{2.1cm}|>{\raggedright\arraybackslash}p{3.6cm}|>{\raggedright\arraybackslash}p{6.6cm}|}
\hline
\textbf{Caso} & \textbf{Risultato} & \textbf{Meccanismo} \\
\hline
\emph{Pegasus} & \emph{``I am a horse.''} (come Java) & MRO: \code{[Pegasus, Horse, Flyer, Mythical, object]} --- vince il \textbf{primo nell'ordine}, non ``la classe'' \\
\hline
\emph{Dragon} & \emph{``I am able to lay eggs.''} (come Java) & MRO: \code{[Dragon, EggLayer, FireBreather, Animal, object]} --- C3 mette \code{EggLayer} prima di \code{Animal} \\
\hline
\emph{FlyingCar} & \textbf{nessun errore!} & prende silenziosamente il primo dell'MRO \\
\hline
\end{tabular}
\end{center}

\wc{Python risolve i conflitti ``come Java''}{
    Coincidenza di risultati, non di meccanismo --- e due differenze
    pesanti, entrambe verificate:
    \begin{enumerate}
        \item \textbf{l'ordine delle basi conta.} Scrivendo
              \code{class Pegasus(Flyer, Horse, Mythical)} invece di
              \code{class Pegasus(Horse, Flyer, Mythical)}, il risultato
              \textbf{cambia}: diventa \emph{``I am able to fly.''}. In
              Java la regola ``istanza batte default'' vale
              \emph{indipendentemente} dall'ordine in cui si scrivono le
              interfacce;
        \item \textbf{il conflitto vero non e' un errore.} Il caso
              \code{FlyingCar} --- due implementazioni indipendenti dello
              stesso metodo --- in Java \textbf{non compila}; in Python
              passa inosservato e vince il primo nell'MRO. E' il caso in
              cui Python e' piu' permissivo e \emph{piu' pericoloso}.
    \end{enumerate}
}

\fonteprove{NL7\_interfacce\_python.py}

\nt{
    Perche' in Python il problema dei \emph{dati} duplicati non si pone
    mai: lo stato di un'istanza vive in \textbf{un solo}
    \code{\_\_dict\_\_}, qualunque sia la profondita' della gerarchia. Non
    esistono sotto-oggetti separati da condividere, quindi non serve nulla
    di analogo all'ereditarieta' virtuale del C++. Resta solo il problema
    dell'\emph{ordine} di risoluzione --- che e' appunto l'MRO.
}

\section{Nested, inner, local: un solo costrutto}

Python non ha alcuna distinzione \emph{sintattica} fra nested, inner,
local e anonymous class: ha un solo costrutto, \code{class}, che si
comporta diversamente a seconda di \textbf{dove} lo si scrive. Le
differenze vanno quindi misurate, non lette nella sintassi.

\section{Misura 1: la classe annidata NON vede lo scope esterno}

\begin{lstlisting}[language=Python]
class Esterna:
    attributo = "ESTERNO"
    class Annidata:
        def prova(self):
            return attributo        # nome NUDO
\end{lstlisting}

Risultato reale: \code{NameError: name 'attributo' is not defined}.

\mprop{Perche'}{
    Il corpo di una classe crea uno \textbf{namespace}, non uno
    \emph{scope} annidato per la ricerca dei nomi: il corpo di
    \code{Annidata} non ``vede'' il corpo di \code{Esterna}. Bisogna
    qualificare: \code{Esterna.attributo}.
    \\[3pt]
    Il comportamento e' quindi lo stesso delle \textbf{nested class del
    C++}: puro \textbf{namespacing}. Python \emph{non} ha inner class.
    Confermato da \code{\_\_qualname\_\_}, che vale
    \code{Esterna.Annidata}.
}

\fonteprove{NL3\_python\_SimpleCalc.py}

\begin{Example}{Un'asimmetria che sorprende}{}
Dentro il corpo di una classe, \emph{mentre viene definita}, i nomi
precedenti \textbf{sono} visibili. Nel file eseguibile:
\begin{lstlisting}[language=Python]
class Calcolatrice:
    class Operazione(ABC): ...
    class Somma(Operazione):        # <- 'Operazione' NUDO funziona qui!
        def applica(self, a, b):
            Calcolatrice.contatore_statico += 1   # <- ma questo va qualificato
            return a + b
\end{lstlisting}
La differenza: \code{Operazione} viene risolto \emph{durante} l'esecuzione
del corpo di \code{Calcolatrice} (quando il namespace locale della classe
e' ancora attivo); \code{contatore\_statico} invece viene risolto
\emph{dopo}, quando il metodo \code{applica} viene chiamato --- e a quel
punto quel namespace non esiste piu'.
\end{Example}

\section{Misura 2: la classe locale SI' cattura (a differenza del C++)}

\begin{lstlisting}[language=Python]
def fabbrica(moltiplicatore):
    class Locale:
        def calcola(self, x):
            return x * moltiplicatore    # variabile LOCALE della funzione
    return Locale
\end{lstlisting}

Risultati reali: \code{fabbrica(3)} produce una classe che moltiplica per
3, \code{fabbrica(5)} per 5.

\mprop{Il contrasto col C++}{
    Questa e' esattamente la riga che in C++ produce
    \emph{``use of local variable with automatic storage from containing
    function''}. In Python funziona perche' i metodi sono
    \textbf{closure}: catturano la \emph{cella} della variabile
    nell'ambiente lessicale che le racchiude.
    \\[3pt]
    Python quindi \textbf{non ha} l'inner class di Java (misura 1), ma
    \textbf{ha} la local class con cattura --- che in C++ richiede una
    lambda.
}

\section{Misura 3: ogni chiamata crea una classe NUOVA}

\begin{lstlisting}
L3, L5 = fabbrica(3), fabbrica(5)
L3 is L5             -> False
isinstance(L3(), L5) -> False
stesso __qualname__  -> True  (fabbrica.<locals>.Locale)
\end{lstlisting}

\mprop{Conseguenza pratica pesante}{
    Sono \textbf{due classi distinte} con lo \emph{stesso nome sorgente}.
    Un'istanza della prima \textbf{non} e' istanza della seconda: quindi
    \code{isinstance} non e' affidabile per confrontare oggetti prodotti
    da chiamate diverse della stessa factory.
    \\[3pt]
    In Java la local class e' compilata \textbf{una volta sola}
    (\code{\$1SottrazioneLimitata}) e le sue istanze condividono il tipo.
    In Python l'oggetto-classe e' \textbf{ricreato a ogni chiamata}, con
    il costo di memoria e tempo che ne segue --- ragione per cui, se non
    serve la cattura per-chiamata, conviene definire la classe fuori.
}

\fonteprove{NL3\_python\_SimpleCalc.py}

\section{Le tre classi dello Strategy in Python}

\begin{lstlisting}[language=Python]
class Calcolatrice:                       # CONTEXT
    contatore_statico = 0                 # attributo di CLASSE

    class Operazione(ABC):                # STRATEGY, annidata
        @abstractmethod
        def applica(self, a, b): ...
        def simbolo(self): return "?"

    class Somma(Operazione):              # FORMA A: annidata
        def applica(self, a, b):
            Calcolatrice.contatore_statico += 1     # va QUALIFICATO
            return a + b
        def simbolo(self): return "+"

    class MoltiplicaPerCampo(Operazione): # FORMA B: "inner", a mano
        def __init__(self, esterno):
            self._esterno = esterno       # enclosing instance ESPLICITA
        def applica(self, a, b):
            return a * b % self._esterno.campo

    def crea_moltiplica_per_campo(self):  # imita esterno.new Inner()
        return Calcolatrice.MoltiplicaPerCampo(self)
\end{lstlisting}

\begin{lstlisting}[language=Python]
# FORMA C: local class, con cattura vera
def crea_sottrazione_limitata(minimo):
    class SottrazioneLimitata(Calcolatrice.Operazione):
        def applica(self, a, b):
            return max(a - b, minimo)     # <- cattura; in C++ e' errore
        def simbolo(self): return f"-(>={minimo})"
    return SottrazioneLimitata()

# FORMA D: la classe davvero "anonima", creata con type()
def divisione_sicura():
    return type(
        "",                                       # nome VUOTO
        (Calcolatrice.Operazione,),
        {"applica": lambda self, a, b: a if b == 0 else a // b,
         "simbolo": lambda self: "/"},
    )()
\end{lstlisting}

Verificato: \code{type(anon).\_\_name\_\_} e' la stringa vuota, e
\code{isinstance(anon, Calcolatrice.Operazione)} e' \code{True} --- e'
una sottoclasse a tutti gli effetti, senza nome.

\nt{
    In Python la forma idiomatica per lo Strategy non e' nessuna di queste
    quattro: e' passare direttamente una \textbf{funzione}. Le classi
    servono quando la strategia ha stato o deve esporre piu' metodi. Il
    documento le mostra tutte perche' l'obiettivo qui e' confrontare
    l'\emph{infrastruttura}, non scegliere lo stile.
}

%==============================================================
\chapter{C: senza classi, senza scope, senza closure}\label{ch:c}
%==============================================================

Il C non ha classi. Le tre categorie delle slide vanno quindi tradotte in
meccanismi diversi --- e nel farlo emergono due limiti che gli altri tre
linguaggi non hanno.

\section{L'equivalente dell'interfaccia: la struct di puntatori}

In C l'``interfaccia'' e' una \code{struct} i cui campi sono
\textbf{puntatori a funzione}: e' letteralmente la vtable, scritta a mano.

\begin{lstlisting}[language=C]
typedef struct {                    /* l'"interfaccia" */
    int (*applica)(const struct Operazione *self, int a, int b);
    const char *simbolo;
} OperazioneVT;

struct Operazione {                 /* la "classe astratta" */
    const OperazioneVT *vptr;
};
\end{lstlisting}

\mprop{Cosa manca, rispetto a un'interfaccia vera}{
    \begin{itemize}
        \item \textbf{nessun controllo di completezza}: se dimentichi di
              assegnare un puntatore, non c'e' alcun errore di
              compilazione --- solo un \code{NULL} che a runtime fa
              crashare. In Java il compilatore \emph{rifiuta} una classe
              che non implementa tutti i metodi;
        \item \textbf{nessun default method}: si puo' lasciare un campo
              a \code{NULL} e far controllare al Context
              (\code{if (vt->f) vt->f(...)}) --- ed e' esattamente cio'
              che fa il VFS del kernel Linux con
              \code{struct file\_operations};
        \item \textbf{nessun conflitto possibile}: non esistendo
              ereditarieta' multipla, il problema che le due regole di
              Java risolvono \emph{non si presenta}. Se servono due
              ``interfacce'', si mettono due vtable nella struct e si
              sceglie a mano quale usare.
    \end{itemize}
}

\fonteprove{NL4\_c\_SimpleCalc.c}

\section{Scoperta 1: annidare una \texttt{struct} NON crea uno scope}

\begin{lstlisting}[language=C]
struct Contenitore {
    int campo;
    struct Annidata {        /* "annidata" solo tipograficamente */
        int x;
    } ann;
};

/* Questa riga COMPILA in C: */
struct Annidata annidata_usata_fuori = { 42 };
\end{lstlisting}

Verificato: il programma compila con \code{gcc -Wall -Wextra} e stampa
\code{42}.

\begin{Prop}{La differenza C / C++}{}
    In \textbf{C} annidare una \code{struct} e' puramente tipografico: il
    tag \code{Annidata} finisce nello \textbf{stesso spazio dei nomi
    globale} del tag esterno. Non c'e' alcun incapsulamento.
    \\[3pt]
    In \textbf{C++} la stessa riga e' un \textbf{errore}:
\begin{lstlisting}
error: variable 'Annidata usata_fuori' has initializer
       but incomplete type
\end{lstlisting}
    perche' li' il nome corretto sarebbe
    \code{Contenitore::Annidata}. E' un caso in cui il C++ \emph{non} e'
    un soprainsieme del C, e mostra che quel poco che le slide
    attribuiscono al C++ (``visibility and naming'') e' comunque piu' di
    quanto offra il C.
\end{Prop}

\fonteprove{NL4\_c\_SimpleCalc.c}

\section{Scoperta 2: le funzioni locali non esistono in C standard}

\begin{lstlisting}[language=C]
int main(void) {
    int fattore = 3;
    int moltiplica(int x) { return x * fattore; }   /* ??? */
    return moltiplica(10);
}
\end{lstlisting}

Compilato con \code{-std=c11 -pedantic}:

\begin{lstlisting}
warning: ISO C forbids nested functions [-Wpedantic]
\end{lstlisting}

Funziona solo come \textbf{estensione GNU} (\code{-std=gnu11}).

\mprop{La conseguenza per lo Strategy}{
    Si puo' definire una \code{struct} dentro una funzione (e quella e'
    davvero locale), ma \textbf{non le si puo' dare un metodo}: le
    funzioni che implementano il comportamento devono stare a livello
    globale. Quindi:
    \begin{itemize}
        \item la \textbf{vtable} e le funzioni-strategia sono sempre
              globali;
        \item la ``localita''' puo' riguardare solo i \emph{dati};
        \item non esistendo closure, il contesto va passato
              \textbf{a mano}.
    \end{itemize}
    Le quattro forme di Java collassano in \textbf{una sola} tecnica.
}

\section{Le tre classi dello Strategy in C}

\subsection{Strategy: una struct di puntatori a funzione}

\begin{lstlisting}[language=C]
struct Operazione;                       /* dichiarazione anticipata */

typedef struct {                         /* la "vtable" a mano */
    int (*applica)(const struct Operazione *self, int a, int b);
    const char *simbolo;
} OperazioneVT;

struct Operazione {                      /* la "classe base astratta" */
    const OperazioneVT *vptr;
};
\end{lstlisting}

\subsection{Context e delega}

\begin{lstlisting}[language=C]
typedef struct {
    int accumulatore;
    const struct Operazione *operazione;
    int campo;                     /* lo "stato dell'oggetto esterno" */
} Calcolatrice;

static void calc_premi_numero(Calcolatrice *c, int n) {
    if (c->operazione == NULL) { c->accumulatore = n; return; }
    /* la delega, identica a strategy.algorithm() di Java/C++ */
    c->accumulatore = c->operazione->vptr->applica(c->operazione,
                                                  c->accumulatore, n);
}
\end{lstlisting}

\subsection{Le tre forme, tradotte}

\begin{lstlisting}[language=C]
/* FORMA A -- l'equivalente della NESTED: nessuno stato proprio */
typedef struct { struct Operazione base; } Somma;

static int somma_applica(const struct Operazione *self, int a, int b) {
    (void)self;                    /* non serve alcuno stato */
    contatore_statico++;           /* usa il "membro static" globale */
    return a + b;
}
static const OperazioneVT SOMMA_VT = { somma_applica, "+" };

/* FORMA B -- l'equivalente dell'INNER: porta il puntatore all'esterno */
typedef struct {
    struct Operazione base;
    const Calcolatrice *esterno;   /* l'"enclosing instance", A MANO */
} MoltiplicaPerCampo;

static int moltiplica_applica(const struct Operazione *self, int a, int b) {
    const MoltiplicaPerCampo *m = (const MoltiplicaPerCampo *)self;
    return a * b % m->esterno->campo;      /* legge lo stato dell'esterno */
}

/* FORMA C -- l'equivalente della LOCAL: il valore "catturato" e' un campo */
typedef struct {
    struct Operazione base;
    int minimo;                    /* la "cattura", esplicita */
} SottrazioneLimitata;
\end{lstlisting}

\mprop{Il ruolo del ``primo membro''}{
    In tutte e tre le forme, \code{struct Operazione base;} e' il
    \textbf{primo} membro. Lo standard C garantisce l'assenza di padding
    all'inizio di una struct, quindi il cast fra
    \code{struct Operazione*} e il tipo concreto e' un \textbf{no-op}: e'
    cio' che permette al Context di trattare uniformemente strategie
    diverse.
}

\subsection{Forma D --- puntatore a funzione e contesto \texttt{void*}}

\begin{lstlisting}[language=C]
typedef int (*FunzioneBinaria)(void *ctx, int a, int b);

typedef struct { int soglia; } CtxMassimo;

static int massimo_con_soglia(void *ctx, int a, int b) {
    CtxMassimo *c = (CtxMassimo *)ctx;
    int m = a > b ? a : b;
    return m > c->soglia ? m : c->soglia;
}

static int applica_funzione(FunzioneBinaria f, void *ctx, int a, int b) {
    return f(ctx, a, b);            /* la delega, senza oggetti */
}
\end{lstlisting}

E' il pattern di \code{qsort\_r}, dei callback POSIX, di GTK e del VFS del
kernel Linux: la coppia (\emph{puntatore a funzione}, \emph{void*}) e'
l'equivalente C di una closure.

\fonteprove{NL4\_c\_SimpleCalc.c}

%==============================================================
\chapter{Confronto finale}\label{ch:confronto}
%==============================================================

\section{La stessa strategia, quattro infrastrutture}

Tutti e quattro i programmi calcolano \code{20 \(\langle\)op\(\rangle\)
6}. Output reali affiancati:

\begin{center}
\renewcommand{\arraystretch}{1.3}
\footnotesize
\begin{tabular}{|>{\raggedright\arraybackslash}p{2.4cm}|>{\raggedright\arraybackslash}p{3.3cm}|>{\raggedright\arraybackslash}p{3.3cm}|>{\raggedright\arraybackslash}p{3.3cm}|}
\hline
\textbf{Forma} & \textbf{Java} & \textbf{C++ / Python} & \textbf{C} \\
\hline
A ``nested'' & \code{static class} & \code{struct} annidata / \code{class} in \code{class} & \code{struct} $+$ vtable \\
\hline
B ``inner'' & \textbf{automatico} (non-static) & puntatore all'esterno \textbf{a mano} & campo \code{esterno} \textbf{a mano} \\
\hline
C ``local'' & con cattura & C++: \textbf{senza} cattura; Python: \textbf{con} & solo dati, metodi globali \\
\hline
D ``anonima'' & \code{new I()\{\dots\}} & lambda / \code{type("")} & \emph{non esiste} \\
\hline
\end{tabular}
\end{center}

\section{Il TIPO della strategia: le quattro dichiarazioni}

\begin{center}
\renewcommand{\arraystretch}{1.35}
\setlength{\tabcolsep}{4pt}
\footnotesize
\begin{tabular}{|>{\raggedright\arraybackslash}p{1.7cm}|>{\raggedright\arraybackslash}p{3.4cm}|>{\raggedright\arraybackslash}p{3.4cm}|>{\raggedright\arraybackslash}p{3.6cm}|}
\hline
 & \textbf{Come si dichiara} & \textbf{Conflitti fra piu' basi} & \textbf{Costo per l'oggetto} \\
\hline
\textbf{Java} & \code{interface} $+$ \code{default}/\code{static} & le \textbf{due regole}; il conflitto vero e' un \textbf{errore} & \textbf{nessuno}: niente stato di istanza \\
\hline
\textbf{C++} & classe astratta pura $+$ dtor virtuale & \textbf{ambiguo}: risolvi a mano, o ereditarieta' \textbf{virtuale} $+$ dominanza & \textbf{un vptr per base} (Pegasus: 24 byte) \\
\hline
\textbf{Python} & \code{ABC} (nominale) o \code{Protocol} (strutturale) & \textbf{MRO}: nessun errore, vince il primo & nessuno: un solo \code{\_\_dict\_\_} \\
\hline
\textbf{C} & \code{struct} di puntatori a funzione & \emph{non esistono}: niente ereditarieta' & un puntatore (il vptr) \\
\hline
\end{tabular}
\end{center}

\mprop{La differenza piu' importante fra i quattro}{
    Solo \textbf{Java} tratta il conflitto fra comportamenti ereditati
    come un \textbf{errore di compilazione} che obbliga il programmatore a
    decidere. C++ lo segnala come ambiguita' (ma solo al punto d'uso, o
    alla definizione della classe se manca il final overrider); Python
    \textbf{non lo segnala affatto} e sceglie in silenzio; il C non puo'
    nemmeno averlo.
}

\section{Le quattro affermazioni da ricordare}

\mprop{Sintesi}{
    \begin{enumerate}
        \item Solo \textbf{Java} ha il legame \emph{implicito} con
              l'istanza esterna (inner class). In C++, Python e C va
              scritto a mano --- e in Java quel campo puntatore esiste
              comunque, solo che non si vede.
        \item Solo \textbf{C++} vieta alla classe locale di leggere le
              variabili automatiche. Per questo la lambda, nel C++
              moderno, e' la vera ``classe locale''.
        \item Solo \textbf{Python} ricrea l'oggetto-classe a ogni chiamata
              della funzione che la contiene, con la conseguenza che
              \code{isinstance} fra classi omonime da chiamate diverse e'
              \code{False}.
        \item Solo il \textbf{C} non ottiene \emph{alcun} incapsulamento
              annidando: il tag resta globale, e senza funzioni locali
              tutte le forme collassano in ``struct $+$ puntatori a
              funzione $+$ contesto esplicito''.
    \end{enumerate}
}

\section{Come scegliere, in pratica}

\begin{center}
\renewcommand{\arraystretch}{1.35}
\footnotesize
\begin{tabular}{|>{\raggedright\arraybackslash}p{5.6cm}|>{\raggedright\arraybackslash}p{6.8cm}|}
\hline
\textbf{Se la strategia\dots} & \textbf{usa\dots} \\
\hline
non ha bisogno dello stato del Context & \textbf{nested/static}: e' la forma di default, non trattiene l'oggetto esterno \\
\hline
deve leggere lo stato \emph{vivo} del Context & Java: \textbf{inner class}; altrove: puntatore/riferimento esplicito \\
\hline
serve solo dentro un metodo, e cattura un parametro & Java: \textbf{local class}; C++: \textbf{lambda}; Python: \textbf{local class}; C: campo nella struct \\
\hline
e' usata una volta sola, in un'espressione & \textbf{anonima} (Java) / \textbf{lambda} (C++, Python) \\
\hline
e' un semplice calcolo senza stato & \textbf{lambda} o funzione libera; in C un puntatore a funzione nudo \\
\hline
\end{tabular}
\end{center}

%==============================================================
\chapter{Esecuzione e verifica}\label{ch:esecuzione}
%==============================================================

\section{I quattro programmi}

\begin{center}
\renewcommand{\arraystretch}{1.3}
\footnotesize
\begin{tabular}{|>{\raggedright\arraybackslash}p{4.5cm}|>{\raggedright\arraybackslash}p{8.0cm}|}
\hline
\textbf{File} & \textbf{Comando} \\
\hline
\code{NL1\_Java\_SimpleCalc.java} & \code{javac NL1\_Java\_SimpleCalc.java \&\& java -cp . NL1\_Java\_SimpleCalc} \\
\hline
\code{NL2\_cpp\_SimpleCalc.cpp} & \code{g++ -std=c++17 -Wall -Wextra NL2\_cpp\_SimpleCalc.cpp -o NL2 \&\& ./NL2} \\
\hline
\code{NL3\_python\_SimpleCalc.py} & \code{python3 NL3\_python\_SimpleCalc.py} \\
\hline
\code{NL4\_c\_SimpleCalc.c} & \code{gcc -std=c11 -Wall -Wextra NL4\_c\_SimpleCalc.c -o NL4 \&\& ./NL4} \\
\hline
\code{NL5\_Java\_Interfacce.java} & \code{javac NL5\_Java\_Interfacce.java \&\& java -cp . NL5\_Java\_Interfacce} \\
\hline
\code{NL6\_interfacce\_}\allowbreak\code{altri\_linguaggi.cpp} & \code{g++ -std=c++17 -Wall -Wextra NL6\_*.cpp -o NL6 \&\& ./NL6} \\
\hline
\code{NL7\_interfacce\_python.py} & \code{python3 NL7\_interfacce\_python.py} \\
\hline
\end{tabular}
\end{center}

\section{Esperimenti: rompili apposta}

\nt{
    Ogni ``ERRORE'' commentato nei sorgenti e' reale: scommentalo e leggi
    il messaggio del compilatore. Sono la parte piu' istruttiva.
    \begin{enumerate}
        \item \textbf{Java} --- crea \code{new MoltiplicaPerCampo()} da
              \code{main} (che e' \code{static}): il compilatore chiede
              un'enclosing instance.
        \item \textbf{Java} --- riassegna la variabile catturata dalla
              local class dopo averla usata: \emph{``local variables
              referenced from an inner class must be final or effectively
              final''}.
        \item \textbf{C++} --- nella local class, sostituisci
              \code{minimoStatico} con \code{minimo}:
              \emph{``use of local variable with automatic storage from
              containing function''}.
        \item \textbf{C++} --- aggiungi \code{static int m;} alla local
              class: \emph{``local class shall not have static data
              member''}.
        \item \textbf{C++} --- nella nested \code{Somma}, prova a leggere
              \code{campoDiIstanza\_}: \emph{``invalid use of non-static
              data member''} --- la prova che non esiste l'inner class.
        \item \textbf{Python} --- nella classe annidata, usa
              \code{contatore\_statico} senza qualificarlo:
              \code{NameError} a runtime (non a ``compile time'': Python
              non ha modo di accorgersene prima).
        \item \textbf{C} --- compila \code{NL4} con \code{-std=c11
              -pedantic} dopo aver aggiunto una funzione dentro una
              funzione: \emph{``ISO C forbids nested functions''}.
        \item \textbf{C} --- rinomina la \code{struct Annidata} in
              \code{Contenitore} e osserva il conflitto: la prova che i
              tag stanno tutti nello stesso spazio globale.
        \item \textbf{Java} --- togli l'override di \code{startEngine} da
              \code{FlyingCar}: \emph{``inherits unrelated defaults''}.
        \item \textbf{Java} --- prova a implementare \code{Bicycle} (4
              metodi astratti) con una lambda: \emph{``is not a functional
              interface''}.
        \item \textbf{C++} --- togli l'override in \code{Pegasus}:
              \emph{``request for member 'identifyMyself' is
              ambiguous''} --- il caso che in Java compila da solo.
        \item \textbf{C++} --- togli \code{virtual} dalle basi di
              \code{EggLayer}/\code{FireBreather}: il caso \code{Dragon}
              diventa ambiguo.
        \item \textbf{Python} --- inverti l'ordine delle basi di
              \code{Pegasus}: il risultato \textbf{cambia} (in Java no).
    \end{enumerate}
}

%==============================================================
\chapter{Fonti}\label{ch:fonti}
%==============================================================

\begin{itemize}
    \item Slide del corso: \emph{L04 --- Focus on the strategy pattern in
          Java} (Giancarlo Succi, Universita' di Bologna), 51 pagine ---
          in particolare le pagg.~6--12 (nested / inner / local class),
          \textbf{13--32 (interfacce, default e static method, le due
          regole di risoluzione dei conflitti, gli esempi Pegasus /
          Dragon / FlyingCar / Mustang / DoIt)}, 33--36 (anonymous class e
          listener), 37--50 (\code{SimpleCalc} con le tre implementazioni
          di \code{ActionListener}).
    \item Bruce Eckel, \emph{Thinking in C++}, Vol.~2, cap.~10
          (\emph{Design Patterns}), sezione \emph{``Strategy: choosing the
          algorithm at runtime''}, e cap.~5 per il policy-based design:
          \url{https://web.mit.edu/merolish/ticpp/TicV2.html\#_Toc53985877}
    \item Wikipedia, \emph{Strategy pattern} --- per la tassonomia dei
          quattro meccanismi (puntatori a funzione, first-class function,
          classi/istanze, reflection):
          \url{https://en.wikipedia.org/wiki/Strategy_pattern}
    \item Cartella \code{prove/strategy/} --- i quattro programmi
          \code{NL1}--\code{NL4} e i test di misura, tutti compilati ed
          eseguiti su questa macchina. \textbf{Ogni output, messaggio
          d'errore e valore di \code{sizeof} citato in questo documento
          proviene da un'esecuzione reale}, non da memoria.
\end{itemize}

\end{document}
